% concatenated from example_paper.tex
%%%%%%%% ICML 2026 EXAMPLE LATEX SUBMISSION FILE %%%%%%%%%%%%%%%%%

\documentclass{article}
% Recommended, but optional, packages for figures and better typesetting:
\usepackage{microtype}
\usepackage{graphicx}
\usepackage{subcaption}
\usepackage{booktabs} % for professional tables
\usepackage{enumitem}

% hyperref makes hyperlinks in the resulting PDF.
% If your build breaks (sometimes temporarily if a hyperlink spans a page)
% please comment out the following usepackage line and replace
% \usepackage{icml2026} with \usepackage[nohyperref]{icml2026} above.
\usepackage{hyperref}
% Attempt to make hyperref and algorithmic work together better:
\newcommand{\theHalgorithm}{\arabic{algorithm}}

% Use the following line for the initial blind version submitted for review:
% \usepackage{icml2026}
% For preprint, use
\usepackage[preprint]{icml2026}

% If accepted, instead use the following line for the camera-ready submission:
% \usepackage[accepted]{icml2026}

\usepackage{amsmath}
\usepackage{amssymb}
\usepackage{mathtools}
\usepackage{amsthm}
\usepackage{bm}
\usepackage{macros}
\usepackage{nicefrac}
\usepackage{thmtools}
\usepackage{thm-restate}
\usepackage[section]{placeins}


% if you use cleveref..
\usepackage[capitalize,noabbrev]{cleveref}

%%%%%%%%%%%%%%%%%%%%%%%%%%%%%%%%
% THEOREMS
%%%%%%%%%%%%%%%%%%%%%%%%%%%%%%%%
\makeatletter
\newcommand*\bigcdot{\mathpalette\bigcdot@{.4}}
\newcommand*\bigcdot@[2]{\mathbin{\vcenter{\hbox{\scalebox{#2}{$\m@th#1\bullet$}}}}}
\makeatother

\newcommand{\mdot}[1]{\overset{\bigcdot}{#1}}
\newcommand{\twodot}[1]{\overset{\bigcdot\bigcdot}{#1}}
\newcommand{\qq}{\mathbf{q}}


\theoremstyle{plain}
\newtheorem{theorem}{Theorem}[section]
\newtheorem{proposition}[theorem]{Proposition}
\newtheorem{lemma}[theorem]{Lemma}
\newtheorem{corollary}[theorem]{Corollary}
\theoremstyle{definition}
\newtheorem{definition}[theorem]{Definition}
\newtheorem{assumption}[theorem]{Assumption}
\theoremstyle{remark}
\newtheorem{remark}[theorem]{Remark}
% \theoremstyle{property}
\newtheorem{property}[theorem]{Property}

% Todonotes is useful during development; simply uncomment the next line
%    and comment out the line below the next line to turn off comments
%\usepackage[disable,textsize=tiny]{todonotes}
\usepackage[textsize=tiny]{todonotes}


\hypersetup{
  breaklinks  = true,
  colorlinks   = true, %Colours links instead of ugly boxes
  urlcolor     = blue, %Colour for external hyperlinks
  linkcolor    = red, %Colour of internal links
  citecolor   = blue %Colour of citations
}



% The \icmltitle you define below is probably too long as a header.
% Therefore, a short form for the running title is supplied here:
\icmltitlerunning{Low-Entropy Outputs of Softmax}

\begin{document}

\twocolumn[
  \icmltitle{Gradient Flow Polarizes Softmax Outputs towards Low-Entropy Solutions}
  % \icmltitle{The Polarizing Dynamics of Gradient Flow on Value-Softmax Model}

% Some implicit biases of softmax parameterization: polarization and low entropy
% Training with softmax parameterization leads to low-entropy solutions
% Polarizing effect of softmax leads to low-entropy solutions 

  % It is OKAY to include author information, even for blind submissions: the
  % style file will automatically remove it for you unless you've provided
  % the [accepted] option to the icml2026 package.

  % List of affiliations: The first argument should be a (short) identifier you
  % will use later to specify author affiliations Academic affiliations
  % should list Department, University, City, Region, Country Industry
  % affiliations should list Company, City, Region, Country

  % You can specify symbols, otherwise they are numbered in order. Ideally, you
  % should not use this facility. Affiliations will be numbered in order of
  % appearance and this is the preferred way.
  \icmlsetsymbol{equal}{*}

  \begin{icmlauthorlist}
      \icmlauthor{Aditya Varre}{yyy}
  \icmlauthor{Mark Rofin}{yyy}
  \icmlauthor{Nicolas Flammarion}{yyy}
    % \icmlauthor{Firstname1 Lastname1}{equal,yyy}
    % \icmlauthor{Firstname2 Lastname2}{equal,yyy,comp}
    % \icmlauthor{Firstname3 Lastname3}{comp}
    % \icmlauthor{Firstname4 Lastname4}{sch}
    % \icmlauthor{Firstname5 Lastname5}{yyy}
    % \icmlauthor{Firstname6 Lastname6}{sch,yyy,comp}
    % \icmlauthor{Firstname7 Lastname7}{comp}
    % %\icmlauthor{}{sch}
    % \icmlauthor{Firstname8 Lastname8}{sch}
    % \icmlauthor{Firstname8 Lastname8}{yyy,comp}
    %\icmlauthor{}{sch}
    %\icmlauthor{}{sch}
  \end{icmlauthorlist}

  \icmlaffiliation{yyy}{Theory of Machine Learning Lab, EPFL}
  % \icmlaffiliation{comp}{Company Name, Location, Country}
  % \icmlaffiliation{sch}{School of ZZZ, Institute of WWW, Location, Country}

  \icmlcorrespondingauthor{Aditya Varre}{aditya.varre@epfl.ch}

  % You may provide any keywords that you find helpful for describing your
  % paper; these are used to populate the "keywords" metadata in the PDF but
  % will not be shown in the document
  \icmlkeywords{Machine Learning, ICML}

  \vskip 0.3in
]

% this must go after the closing bracket ] following \twocolumn[ ...

% This command actually creates the footnote in the first column listing the
% affiliations and the copyright notice. The command takes one argument, which
% is text to display at the start of the footnote. The \icmlEqualContribution
% command is standard text for equal contribution. Remove it (just {}) if you
% do not need this facility.

% Use ONE of the following lines. DO NOT remove the command.
% If you have no special notice, KEEP empty braces:
\printAffiliationsAndNotice{}  % no special notice (required even if empty)
% Or, if applicable, use the standard equal contribution text:
% \printAffiliationsAndNotice{\icmlEqualContribution}

% \begin{abstract}
%   Softmax has been integral part of transformers self-attention blockas. In this article, we look at some dynamical properties of training with softmax. To this extent, we consider the loss function of form $\bloss\left( \VV \sf(\aa) \right)$ where $\VV$ is a trainable value matrix and $\aa$ is a trainable score vector. We analyze the gradient flow dynamics of this model and show that the softmax parameterization leads to solutions where the softmax output has low entropy. We demonstrate the generality of our results by considering various loss functions such as logistic loss, regression etc.. Finally, we discuss some implications of our results to empirical observations particularly attention sinks.
% \end{abstract}

% \begin{abstract}
%   The softmax function is a integral component of the self-attention mechanism in Transformers.  This article investigates the dynamical properties associated with training softmax-based models. Specifically, we analyze the gradient flow dynamics of a loss function defined as $\ell(\mathbf{V} \sf(\mathbf{a}))$, where $\mathbf{V}, \mathbf{a}$ respectively are learnable value matrix, attention vector and $\sf$ is a softmax function. Our analysis reveals that Softmax parameterization inherently drives the optimization toward solutions characterized by low-entropy outputs. We demonstrate the universality of these findings across various objectives, including logistic loss and regression. Furthermore, we discuss the practical implications of these theoretical results, offering new insights into empirical phenomena such as "attention sinks."
% \end{abstract}

% \begin{abstract}
% The softmax function is an integral component of the self-attention mechanism in Transformers. This article investigates the dynamical properties associated with training softmax-based models. Specifically, we analyze the gradient flow dynamics of a loss function defined as $\ell(\mathbf{V} \sigma(\mathbf{a}))$, where $\mathbf{V}$ and $\mathbf{a}$ are a learnable value matrix and an attention vector, respectively, and $\sigma$ denotes the softmax function. Our analysis reveals that gradient flow over this softmax parameterization inherently drives the optimization toward solutions characterized by low-entropy outputs. We demonstrate the universality of these findings across various objectives, including logistic loss and regression. Furthermore, we discuss the practical implications of these theoretical results, offering new insights into empirical phenomena such as \emph{attention sinks}.
% \end{abstract}

% \begin{abstract}
%   The softmax function is an integral component of the self-attention mechanism in Transformers. This article investigates the dynamical properties associated with training softmax-based models. Specifically, we analyze the gradient flow dynamics of a loss function defined as $\bloss(\mathbf{V} \sigma(\mathbf{a}))$, where $\mathbf{V}$ and $\mathbf{a}$ are a learnable value matrix and learnable attention vector, respectively, and $\sigma$ denotes the softmax function. Our analysis reveals that gradient flow over this softmax parameterization inherently drives the optimization toward solutions characterized by low-entropy outputs. We demonstrate the universality of these findings across various objectives, including logistic loss and square loss. Furthermore, we discuss the practical implications of these theoretical results, offering new insights into empirical phenomena such as \emph{attention sinks}.
% \end{abstract}

% \begin{abstract}
%   % Understanding the dynamics of training self-attention models is challenging yet crucial for advancing deep learning.
%   In this article, we analyze the gradient flow dynamics of a loss function over a \emph{value-softmax} model defined as $\bloss(\mathbf{V} \sigma(\mathbf{a}))$, where $\mathbf{V}$ and $\mathbf{a}$ are a learnable value matrix and learnable attention vector, respectively, and $\sigma$ denotes the softmax function. Note that such a product of a matrix and a softmax vector is ubiquitous in self-attention mechanisms in Transformers.   Our analysis reveals that gradient flow over this value-softmax parameterization inherently drives the optimization toward solutions characterized by low-entropy outputs. We demonstrate the universality of these findings across various objectives, including logistic loss and square loss. Furthermore, we discuss the practical implications of these theoretical results, offering new insights into empirical phenomena such as \emph{attention sinks, massive activations}.
% \end{abstract}

\begin{abstract}
Understanding the intricate non-convex training dynamics of softmax-based models is crucial for explaining the empirical success of transformers. In this article, we analyze the gradient flow dynamics of the \emph{value-softmax} model, defined as $\mathcal{L}(\mathbf{V} \sigma(\mathbf{a}))$, where $\mathbf{V}$ and $\mathbf{a}$ are a learnable value matrix and attention vector, respectively. As the \emph{matrix times softmax vector} parameterization constitutes the core building block of self-attention, our analysis provides direct insight into transformer's training dynamics. We reveal that gradient flow on this structure inherently drives the optimization toward solutions characterized by low-entropy outputs. We demonstrate the universality of this polarizing effect across various objectives, including logistic and square loss. Furthermore, we discuss the practical implications of these theoretical results, offering a formal mechanism for empirical phenomena such as \emph{attention sinks} and \emph{massive activations}.
\end{abstract}

\section{Introduction}

Despite the widespread adoption of large language models (LLMs), our understanding of the internal mechanisms by which they process information remains limited. In particular, the self-attention mechanism and the representations it induces play a central role in model behavior, yet many aspects of their functioning remain in the early stages of understanding. Recent work has begun to shed light on this opacity by identifying recurring learned structures within transformer architectures, such as induction heads~\citep{olsson2022context} and attention sinks~\citep{xiao2024efficient}. These components have attracted significant attention as crucial building blocks underlying complex model capabilities.

Building on these efforts, identified components are often examined in isolation to assess how architectural design choices, training algorithms and hyperparameters, as well as properties of the training data, contribute to their emergence. A theoretical understanding of these factors is a necessary step toward improving the robustness, efficiency, and safety of large language models. 
%In this work, we pursue this objective by focusing on one such crucial component of transformer architectures: low-entropy attention patterns, and question if there are solely implied by the task. 
%In this work, we pursue this objective by studying a broader phenomenon that underlies many of these components: \emph{low-entropy (sparse) attention patterns}, where the attention distribution concentrates most of its mass on a small number of tokens. We ask whether such sparsification is a functional consequence of the task, or instead reflects an implicit preference of optimization and parameterization.
In this work, we pursue this objective by studying a broader phenomenon that underlies many of these components: \emph{low-entropy (sparse) attention patterns}, where the attention distribution concentrates most of its mass on a small number of tokens, and by asking whether this sparsification is solely enforced by the task.

% In the body:
\begin{figure}[t]
  \centering
  % Tight spacing between subfigures (optional)
  \setlength{\tabcolsep}{2pt}

  \begin{subfigure}[t]{0.31\columnwidth}
    \centering
    \includegraphics[width=\linewidth]{figs/attention_map_softmax_0_1_0.pdf}
    \caption{Softmax}
    \label{fig:attn-softmax}
  \end{subfigure}\hfill
  \begin{subfigure}[t]{0.31\columnwidth}
    \centering
    \includegraphics[width=\linewidth]{figs/attention_map_sigmoid_0_1_0.pdf}
    \caption{Sigmoid}
    \label{fig:attn-sigmoid}
  \end{subfigure}\hfill
  \begin{subfigure}[t]{0.31\columnwidth}
    \centering
    \includegraphics[width=\linewidth]{figs/attention_map_linear_0_1_0.pdf}
    \caption{Linear}
    \label{fig:attn-linear}
  \end{subfigure}

  \caption{Representative attention patterns of the 2nd Transformer layer solving an induction task. Sigmoid and linear attentions are trained without additional normalization (see Table \ref{table:attention-types}). An attention sink clearly emerges for softmax but not other attentions.}
  \label{fig1}
  \vspace{-10pt}
\end{figure}

%As a motivating example, we consider \emph{attention sinks}~\citep{xiao2024efficient}, sparse attention patterns characterized by attention mass that is highly concentrated on a specific token, often the first token in the sequence. Subsequent studies~\citep{sun2024massive, yu2024unveiling} report that sink-like patterns can also emerge at later positions, including tokens with limited semantic relevance and even input-agnostic locations. %These findings suggest that such sparse attention patterns may not be explained solely by task semantics or sequence structure.
% common hypothesis is that low-entropy attention is useful because concentrating mass on a small set of tokens induces a systematic bias in the attention output. These findings nonetheless suggest that such sparse attention patterns may not be explained solely by task semantics or sequence structure.
%More broadly, they motivate the following question: 
%\emph{Why does attention become sparse in the first place, and does low-entropy attention reflect a functional requirement of the task or an implicit preference induced by optimization and parameterization?}

%These findings suggest that sinks may not be explained solely by task semantics or sequence structure. A common hypothesis is that sinks are useful because they induce explicit biases in the attention output, but it remains unclear \emph{why} sparse attention should arise in the first place, and whether it reflects a functional requirement of the task or an implicit preference of optimization and parameterization.


As a motivating example, we consider \emph{attention sinks}~\citep{xiao2024efficient}: sparse attention patterns characterized by attention mass that is highly concentrated on one token, often the first token in the sequence. A common hypothesis is that such low-entropy patterns are useful because concentrating mass on a small set of tokens induces a systematic bias in the attention output; however, the need for such bias is not always clear from task semantics or sequence structure. Moreover, some softmax-free attention types do not form attention sinks (Figure \ref{fig1}), pointing to the role of architecture in modulating their emergence.
This motivates us to ask the following question: 
\emph{Why does attention become sparse in the first place, and does low-entropy attention reflect a functional requirement of the task or an implicit preference induced by optimization and parameterization?}


We investigate this question by isolating the minimal computation underlying a single attention head. The output of a head can be written as a value matrix multiplied by a softmax attention distribution, i.e., $\VV\sf(\aa)$. In this form, an attention sink corresponds to $\sf(\aa)$ becoming low-entropy (near one-hot), concentrating mass on a single token. We therefore study a simplified \emph{value--softmax} model $\VV\sf(\aa)$, where both $\VV$ and $\aa$ are trainable, and analyze whether gradient-based training exhibits an \emph{implicit bias toward sparsification} even when many dense solutions exist. Concretely, we consider several losses with a fixed target vector, analyze gradient-flow dynamics, and complement the theory with empirical evidence.


% Transformers are important 
% Despite the widespread use of Large Language models, the understanding of the core mechanisms with which they process information is still in the nascent stages.  Self-attention mechanism and its internal representation play a crucial role in this information processing. 

% Contributing to this line of research of understanding, some components of the transformer like induction heads, attention sinks have been identified and have received much attention lately.

% Once such components are identified, they are then studied in isolation to unveil what parts of architecture, training algorithm and its hyper-parameters, properties of data influence the formation of these structures.  Such theoretical foundations are essential to make them robust, efficient and safe. 


% Once identified, such components are typically studied in isolation to determine how architectural choices, training algorithms, and hyperparameters, as well as data properties, contribute to their emergence. Developing a theoretical understanding of these factors is essential for improving the robustness, efficiency, and safety of large language models.





%\emph{why} sparse attention should arise in the first place, and whether it reflects a functional requirement of the task or an implicit preference of optimization and parameterization.

%, such as attention sinks.
%attention sinks. 
% ~\citet{xiao2024efficient} identified the attention sinks which are sparse activation patterns where the attention is localized to first token. This phenomenon of localization is leveraged extensively for resource optimization like KV cache optimization, long context generation, quantization and pruning. Due to this widespread applications, attention sinks have received a lot of interested.  

% Although initial work has shown that the sinks are  formed at the first token, \citep{sun2024massive, yu2024unveiling} have identified that the sinks can emerge at later tokens of less semantic importance and also input agnostic position.  One of their main usefulness of these patterns is to create certain \textit{explicit biases}. Why such biases are created through only sparse activation patterns ? 
% % There are many attention parameters which are spiky and of low entropy - either you call it sparse attention patterns or attention sinks or . 
% This is the motivation of this work. To establish theoretical foundation for this phenomenon through a simplified model and extend implications to more real-world-like architectures. 

%Attention sinks~\citep{xiao2024efficient} are such sparse attention patterns characterized by the attention mass that is highly concentrated on a specific token, often the first token in the sequence. 
%This phenomenon is broader and more puzzling than early observations suggested. Subsequent 
%Studies~\citep{sun2024massive, yu2024unveiling} report that sink-like patterns can also emerge at later positions, including tokens with limited semantic relevance and even input-agnostic locations. 
%
%These findings suggest that these sparse attention patterns may not be explained solely by task semantics or sequence structure. A common hypothesis is that sinks are useful because they induce explicit biases in the attention output,  but it remains unclear \emph{why} sparse attention should arise in the first place, and whether it reflects a functional requirement of the task or an implicit preference of optimization and parameterization.

%However, subsequent studies~\citep{sun2024massive, yu2024unveiling} suggest that attention sinks are not merely a BOS-specific curiosity. Sink-like patterns can emerge at later positions, including tokens with limited semantic relevance and even input-agnostic locations. Rather than indicating a semantically meaningful routing decision, these observations are consistent with a different interpretation: sinks may serve as a \emph{content-independent anchor} that is available for every input and can induce an explicit bias in the attention output. What remains unclear is why such an anchor is realized specifically through \emph{sparse} attention---i.e., why gradient-based training tends to concentrate probability mass onto a single token (or a small set of tokens), as opposed to distributing this effect across many tokens or implementing the bias through other components of the architecture.

%While early work primarily observed attention sinks at the first token, subsequent studies~\citep{sun2024massive, yu2024unveiling} have demonstrated that such patterns can also emerge at later positions, including tokens with limited semantic relevance and even input-agnostic positions. It is hypothesized that one of the key utilities of these attention patterns lies in their ability to induce explicit biases. However, it remains unclear why such biases arise specifically from sparse attention activations.

%We investigate this question by isolating the minimal computation underlying a single attention head. The output of a head can be written as a value matrix multiplied by a softmax attention distribution, i.e., $V\,\sf(a)$. In this form, an attention sink corresponds to $\sf(a)$ becoming low-entropy (near one-hot), concentrating mass on a single token. We therefore study a simplified \emph{value--softmax} model $V\sf(a)$, where both $V$ and $a$ are trainable, and analyze whether gradient-based training exhibits an \emph{implicit bias toward sparsification} even when many dense solutions exist. Concretely, we consider several losses with a fixed target vector, analyze gradient-flow dynamics, and complement the theory with empirical evidence.
%In particular,  we focus on  a value-softmax model $V\sf(\aa)$, i.e., the model is product of a matrix $\VV$ and softmax of a vector $\aa$, where both are trainable. Note that product of matrix and softmax function is ubiquitous in the self-attention based architectures. The objective of this model is to create an implicit attention bias term~\citep{sun2024massive}; to this end, we study various loss functions with a fixed target vector. By studying the dynamics of gradient flow on this model, complemented by explicit empirical evidence, we make the following contributions:

Our contributions are:
\begin{itemize}[noitemsep, topsep=1pt]
\item We study the value-softmax model with logistic loss and show that among many possible solutions, the gradient flow is biased towards the sparse solution and the output of the softmax, i.e., $\sf(\aa)$ converges to a \emph{one-hot} vector. This sparsification is attributed to the polarizing dynamics of the gradient flow.
%
\item We extend these polarizing dynamics to the case of regression and attribute the sparsification to the speed of convergence. We further discuss how alternative nonlinearities and normalization schemes modulate (or prevent) this polarization.
%
%\item We present the implications of our results for attention sinks, we show that attention sinks and massive activations can happen as a result of implicit bias of the softmax parameterization. 
\item We connect these results to transformer attention, arguing that attention sinks and massive activations can emerge as a consequence of the implicit bias induced by the softmax parameterization, helping explain sink formation.% even at semantically uninformative or input-agnostic positions.
%
\item We provide empirical evidence on how normalization, both its presence and absence can impact the proportion of attention sinks specifically in the formation of induction heads. We also discuss the pitfalls of such sparsification bias, where a dispropotionate burden of decision making is attributed to a single token.
\end{itemize}


% What are the applications?
% \begin{itemize}
%   \item realtion with attention sinks.
%   \item solutions with low entropy with softmax. 
% \end{itemize}
\subsection{Related Work}

Our work lies at the intersection of mechanistic studies of attention in transformers, the emergence of low-entropy attention patterns, optimization-driven implicit bias in softmax-parameterized models, and polarization-style dynamics in gradient-flow.

\paragraph{Attention sinks and massive activations.}
\citet{xiao2024efficient} identified that large proportional of attention scores is focused on the first token and termed such tokens as attention sinks.~\citet{yu2024unveiling,sun2024massive} also identified that the attention sinks do not necessarily correspond to the first token. Explaining the phenomenon,
\citet{sun2024massive} attribute the sink phenomenon to the massive activations which function as implicit attention biases. \citet{xiao2024efficient,gu2025attention} attribute the sink phenomenon to softmax normalization. \citet{gu2025attention} also studies the impact of data distribution, optimization, loss function and architecture on the formation of attention sinks. More recently, \citet{zhang2025attention} describe %attention sinks catch all the embeddings and tag them with a fixed embeddings which is later used to cluster these tokens in subsequent layers. 
attention sinks as a mechanism intended to cluster the hidden representations of input tokens, and \citet{barbero2025llms} conclude that they are needed to prevent representational collapse.
\citet{guo2024active} study the formation of attention sinks using a simple Markovian data model and identify that heads may attend to sinks for some but not all input domains.
%become sinks for specific tokens and not others. %For vision tasks, \citet{darcet2024vision} provide additional tokens (registers) to act as sinks.

Substantial line of research has employed attention sinks to achieve practical benefits, particularly improvements in
%This localization property has been extensively leveraged for resource optimization techniques such as 
key–value cache compression~\citep{guo-etal-2024-attention}, long-context generation~\citep{xiao2024efficient,Yang_2025_ICCV}, quantization~\citep{son-etal-2024-prefixing}, and feature map smoothness~\citep{darcet2024vision}.
%Owing to these practical benefits, attention sinks have attracted significant research interest.

\paragraph{Softmax alternatives and entropy collapse.}
% Low entropy in attention distributions can lead to representational collapse or ``over-squashing'' . This variance sensitivity is a known cause of entropy collapse .  proposed methods to stabilize training by preventing this collapse.

Numerous alternatives to softmax attention have been proposed, mostly from the perspective of improved computational complexity and parallelizability, including sparsemax \citep{martins2016softmax}, linear attention \citep{pmlr-v119-katharopoulos20a, qin2022cosformer, wortsman2023replacing}, sigmoid attention \citep{ramapuram2025theory}, and others.

It has been shown that softmax leads to much lower attention entropy than its linear alternatives \citep{hong-lee-2025-variance}. However, the impact of this difference is not entirely clear: \citet{stabilizing-transformer-training} argue that entropy collapse leads to training instability and attempt to prevent it from happening, \citet{zhang2024hedgehog} state that low-entropy attention patterns are useful and design an attention operation to preserve them.

To the end of explaining entropy collapse, \citet{deng2024sparse} show that it manifests with high sequence length due to the log-scale increase in logits. Further, \citet{bao2024self} tie attention localization to the eigenspectrum of QK matrices. 

\paragraph{Sparse attention patterns.} Independently from the works on entropy collapse, it has been shown that softmax Transformers often implement sparse solutions, driving a lot of progress in Transformer interpretability \citep{olsson2022context, lindsey2025biology}. \citet{singh2024needs, zucchet2025emergence, yuksel2025incremental, varre2025learning} study the conditions required for the emergence of sparse attention patterns. 
\citet{nanda2023factfinding} show that LLMs exhibit subcircuits for factual recall, and
\citet{zucchet2025language} analyze the evolution dynamics of attention for this task.

\paragraph{Polarization and population dynamics.} From a technical perspective, the dynamical system analysis of our value-softmax model shares significant structural similarities with replicator dynamics in evolutionary game theory~\citep{taylor1978evolutionary, smith1982evolution}; see \citet{sorin2020replicator} for a survey and opinion polarization on social networks~\citep{proskurnikov2015opinion}. In these frameworks, the evolution of a state $x \in \mathbb{R}^{p}$ is generally governed by dynamics of the form:\begin{align*}\dot{x}_i \propto ( f(x_i) - \bar{f}(x) ),\end{align*}where $\bar{f}(x)$ represents the population (weighted-)aggregate. Because the rate of change is driven by the deviation from this aggregate, parameters tend to diverge from the mean. This mechanism underpins phenomena such as \emph{survival of the fittest} in biological contexts and the hardening of polarized opinions in social dynamics. 


%\citet{hong-lee-2025-variance}
%
%Several alternatives to the standard softmax have been proposed to mitigate these issues:
%\begin{itemize}
%    \item \textbf{Sigmoid Attention:} Analyzed by .
%    \item \textbf{ReLU Attention:}  explore replacing softmax with ReLU in vision transformers.
%    \item \textbf{Sparsemax:}  proposed a sparse model of attention that maps the Euclidean projection onto the simplex.
%\end{itemize}

% \paragraph{Attention Sinks}
% \begin{itemize}
%     \item \href{https://arxiv.org/pdf/2504.02732}{Why do LLMs attend to the first token?}
%     \item \href{https://arxiv.org/pdf/2502.00919}{Attention Sinks: A `Catch, Tag, Release' Mechanism for Embeddings} \citep{zhang2025attention}
%     \item \textbf{Xiao et al. (2023):} ``Efficient Streaming Language Models with Attention Sinks.''
%     \begin{itemize}
%         \item \textit{Note:} They empirically observed that softmax forces the model to dump attention on a ``sink'' token (often the first token) when no other token is relevant. (Connection: Proof of ``low entropy solutions'' provides theoretical justification).
%     \end{itemize}
%     \item \textbf{Gu et al. (2024):} \href{https://arxiv.org/pdf/2402.17762#page=2.99}{Massive Activations}
%     \begin{itemize}
%         \item \textit{Note:} Discusses how ``outliers'' in feature dimensions relate to polarized $a$ vectors.
%     \end{itemize}
% \end{itemize}

%\paragraph{Low Entropy of Attention}



%\begin{itemize}
%    \item \href{https://arxiv.org/pdf/2406.04267}{Transformers need glasses! Information over-squashing in language tasks} (Note: Representation %collapse may be not related).
%    \item \href{https://aclanthology.org/2025.emnlp-main.421v2.pdf}{Variance Sensitivity Induces Attention Entropy Collapse in Transformers}
%    \item \href{https://machinelearning.apple.com/research/stabilizing-transformer-training}{Stabilizing Transformer Training by Preventing Attention %Entropy Collapse}
 %   \item \href{https://arxiv.org/pdf/2402.02098}{arXiv:2402.02098} (General reference regarding attention entropy).
%\end{itemize}

%\paragraph{Alternatives to Softmax}
%1. quadratic complexity of standard softmax attention motivates the research of alternative more efficient architectures.
%\begin{itemize}
%    \item \href{https://proceedings.mlr.press/v119/katharopoulos20a.html?ref=mackenziemorehead.com}{Transformers are RNNs: Fast Autoregressive %Transformers with Linear Attention}
%    \item \href{https://arxiv.org/abs/2202.08791}{cosFormer: Rethinking Softmax in Attention}
%    \item \href{https://arxiv.org/abs/2402.04347}{The Hedgehog \& the Porcupine: Expressive Linear Attentions with Softmax Mimicry} -- these guys also %measure entropy of different attention architectures
%    \item \href{https://arxiv.org/pdf/2409.04431}{Theory, Analysis, and Best Practices for Sigmoid Self-Attention}
%    \item \href{https://arxiv.org/pdf/2309.08586}{Replacing softmax with ReLU in Vision Transformers}
%    \item \href{https://proceedings.mlr.press/v48/martins16.pdf}{From Softmax to Sparsemax: A Sparse Model of Attention and Multi-Label Classification}
%\end{itemize}

% \paragraph{Polarization: Population Dynamics and Consensus}
% \begin{itemize}
%     \item \textbf{Opinion Dynamics:} \href{https://arxiv.org/pdf/1508.05034}{Opinion Dynamics in Social Networks with Hostile Camps: Consensus vs. Polarization}
%     \item \textbf{Replicator Dynamics:} \href{https://perso.imj-prg.fr/sylvain-sorin/wp-content/uploads/sites/11/2023/10/20.RD_.JDG_.pdf}{Replicator dynamics: OLD and new}
%     \item \textbf{Cultural Polarization:} Small Worlds and Cultural Polarization
%     \item \textbf{Softmax in Game Theory:} \href{https://arxiv.org/pdf/1704.00805}{On the Properties of the Softmax Function with Application in Game Theory and Reinforcement Learning}
%     \item \textbf{ESS:} \href{http://www.dklevine.com/archive/refs4457.pdf}{Evolutionarily Stable Strategies and Game Dynamics}
%     \item \textbf{Community Ecology:} \href{https://stefanoallesina.github.io/Theoretical_Community_Ecology/game-theory-and-replicator-dynamics.html}{Game theory and replicator dynamics}
% \end{itemize}
% \citep{vivien2022label}


\section{Problem Setup}

\paragraph{Notations.} For $n \in \mathbb{N}$, let $[n] := \{0,1,\ldots,n{-}1\}$. Bold letters denote quantities of interest which includes time-dependent variables and their limiting values. For a vector $x \in \mathbb{R}^d$, $x_i$ denotes its $i^{th}$ coordinate. For a matrix $X \in \mathbb{R}^{m \times n}$, $X_i$ denotes its $i^{th}$ column. The notation $\|\cdot\|$ denotes the Euclidean norm for vectors and the Frobenius norm for matrices. For a vector $a \in \mathbb{R}^d$, $\mathrm{diag}(a)$ denotes the diagonal matrix with entries given by $a$. The all-ones vector is denoted by $\one$ and its dimension is inferred from the context. For vectors $u, v \in \mathbb{R}^d$, $u \odot v$ denotes the Hadamard (elementwise) product. 

The softmax function $\mathrm{\sf} : \mathbb{R}^p \to \mathbb{R}^p$ is defined by
\begin{align*}
  \sf(a)_i = \frac{\exp(a_i)}{\sum_{j=1}^{p} \exp(a_j)}.  
\end{align*}
We first present a single self-attention block and then simplify it to a model we call the value–softmax model.

\paragraph{Self attention.} 
%
Let $X \in \R^{T \times d}$ be a sequence of input token embeddings and let $Q,K,V$ denote the query, key, and value matrices of the attention block. The updated representation after the attention block is given by 
\begin{align*}
  X_{+} = \sf\left( X QK^{\top} X^{\top} \right) X V^{\top}. 
\end{align*}
Let $z$ denote the embedding of the last input token. Its updated embedding is 
\begin{align*}
  z_{+}  = V X^{\top} \sf\left( X K Q^{\top} z\right) = \bar{V}^{\top} \sf\left( \bar{a} \right),
\end{align*}
where we define $\bar{V} = XV^{\top} \in \R^{T \times d}$ and $\bar{a} = X K Q^{\top} z \in \R^{T}$. 
%Recall the definition of softmax function $\sf: \R^{T} \to \R^{T}$ is given by\begin{align*} \sf(a)_i = \frac{\exp(a_i)}{\sum_{j=1}^{T} \exp(a_j)}.\end{align*}
Such an update has the following structure: a trainable matrix $\bar{V}$ multiplied by softmax of a vector $\bar{a}$. The embedding 
\begin{align*}
z_{+} = \sum_{i} v_i \sigma(a)_i 
\end{align*}
can be represented in multiple ways. For example, the scores $\sigma_{i}'s$ may be uniform, in which case the vectors  $v_i'$ average to $z_{+}$; alternatively, the scores may be sparse(e.g., one-hot) in which case a single vector $v_i$ is selected to be $z_{+}$. 

We would like to understand which of these regimes the optimization procedure tends to favor. Note that in the above derivation we ignore the contribution of the value matrix $V$ to the updates of the other tokens and assume independence of the $z_{+}$. 
While these assumptions do not hold in general, they motivate a simplified model, defined next, that provides insight into the behavior induced by the softmax parameterization.

\paragraph{A value-softmax model.}
%
Let $\VV \in \R^{p \times p}$ and $\aa \in \R^{p}$ denote the trainable value matrix and score vector, respectively.  Let $\LL{\VV, \aa}$ be a loss function that depends only on the product  $\bbeta = \VV \sf(\aa)$:
\begin{align}
  \LL{\VV, \aa} = \bloss\left( \VV \sf(\aa) \right) = \bloss\left( \bbeta \right).
\end{align} 
The loss does not depend on $\aa$ and $\VV$ separately but only through the combination $\VV \sf(\aa)$. This model is a stylized simplification of the self-attention block.
%
Re-parameterized models are a popular setting to gain insight into the training dynamics of non-convex optimization problem~\citep{lireparam2022}.
%
For convenience, we denote $\sos = \sf(\aa)$.  

\paragraph{Gradient flow.}
%
We study the dynamics of the continuous time limit of gradient descent, the gradient flow. Although this analysis does not capture discrete-time effects, adaptive step sizes, or stochasticity, it nevertheless captures the principal behavior of first-order optimization and provides useful insight into the training process. 

The gradient flow for the parameters $\VV$ and $\aa$ is given by the differential equations 
\begin{align} \label{eq:gf-L}
  \der{\VV} &= - \nabla_{\VV} \LL{\VV, \aa} \der{t},  \,\,
  \der{\aa} = - \nabla_{\aa} \LL{\VV, \aa} \der{t}.
\end{align}
We exploit the  reparameterization $\LL{\VV,\aa} = \bloss\left( \VV \sf(\aa)\right)$ to obtain explicit expressions for the gradients. Using the chain rule on $\bbeta = \VV \sf\left( \aa \right)$, the gradient flow can be written as 
\begin{align}
  \der{\VV}  &= - \nabla_{\bbeta} \bloss(\bbeta) \sos^{\top} \der{t}, \\
  \der{\aa}  &= - \left[ \diag{\sos} - \sos \sos^{\top} \right] \VV^{\top} \nabla_{\bbeta} \bloss(\bbeta).
\end{align}
The Jacobian of the softmax $\diag{\sos} - \sos \sos^{\top}$ plays a central role in shaping the dynamics of  $\aa$.

\section{Polarizing effect of softmax}
%
In this section, we present the main results for the value–softmax model. We begin with the logistic loss and then discuss regression losses as well as other nonlinearities and normalization functions.

\subsection{Logistic loss}
%

% Recall that the goal is to capture a scenario where an attention head creates an implicit attention bias. The value-softmax model captures the head and say it needs to act as a specific attention bias say $\bbeta_*$. Coming to the loss function, the logistic loss is motivated by the downstream propagation of attention biases in transformers. When an attention head generates an implicit attention bias,  this bias is consumed by subsequent attention layers where it influences binary decisions~\citep{zhang2025attention}. 

% Recall that the goal is to capture a scenario where an attention head creates an implicit attention bias. The value-softmax model captures this head, which needs to produce a specific attention bias $\bbeta_*$. For the loss function, we choose the logistic loss, motivated by the downstream propagation of attention biases in transformers. When a head generates an implicit attention bias, this bias is consumed by subsequent attention layers where it influences discrete decisions~\citep{zhang2025attention}, or directly by cross-entropy loss (which reduces to logistic loss in the binary case).


We consider a binary classification setting, mirroring standard transformer training where the final representation is optimized with cross-entropy (i.e., logistic loss in the binary case). 
%
To isolate the effect of the value–softmax parameterization, we consider a stylized teacher setting in which the desired decision rule is given by a fixed linear classifier $\bbeta_* \in \R^{p}$.
%
The attention block produces $\bbeta$, and the logistic loss encourages a large positive margin $\scal{\bbeta_*} {\bbeta}$, i.e., agreement with the classifier induced by $\bbeta_*$.

This setting leads to the following formulation of the loss as a function of $\VV,\aa$:
\begin{align} \label{eq:logistic-loss-V-a}
  \LL{\VV, \aa} = \log\left( 1 + \exp \scal{\bbeta_*} {\, - \VV \sf(\aa)} \, \right),
\end{align}
%for the target vector $\bbeta_* \in \R^{p}$. 
The gradient flow dynamics are given by 
\begin{align*}
  \der{\VV} &= - \gamma(\bbeta) \, \bbeta_* \sos^{\top} \der{t}, \\
  \der{\aa} &= - \gamma(\bbeta) \, \left[ \diag{\sos} - \sos \sos^{\top} \right] \VV^{\top} \bbeta_* \der{t}.
\end{align*}
where $\gamma(\bbeta) = \left(1 + \exp\scal{\bbeta_*}{\bbeta}\right)^{-1}$.
Letting  $\uu = \VV^{\top} \bbeta_*$, the dynamics reduced to a couples system over two vectors $\aa, \uu \in \mathbb{R}^p$ 
\begin{align} \label{eq:logistic-loss-u-a}
  \mdot{\uu} &= - \gamma(\bbeta) \, \nor{\bbeta_*}^2 \sos, \\
  \mdot{\aa} &= - \gamma(\bbeta) \, \left[ \diag{\sos} - \sos \sos^{\top} \right] \uu.
\end{align}
%

\paragraph{Intuition from replicator dynamics.}
%
Inspecting the gradient-flow equations, the evolution of the scores is driven  by the centered signal
%Looking at the time derivative of $\aa$ and second derivative of $\uu$, 
\begin{align*}
\mdot{\aa} \sim \left( \uu  -  \scal{\uu}{\sos} \one  \right) , \quad \twodot{\uu} \sim  \left( \uu  -  \scal{\uu}{\sos} \one  \right),
\end{align*}
where we omit scaling factors and normalization terms for clarity.  
%This differential equality holds similarity to replicator dynamics in population dynamics and evolutionary game theory and dynamics of consensus and polarization of opinions in social networks. 

Since $\sos=\sf(\aa)$ lies on the simplex, $\langle \uu,\sos\rangle$ is the $\sos$-weighted average of the coordinates of $\uu$. Thus, the term $\uu - \langle \uu,\sos\rangle \one$ measures each coordinate's deviation from the average.
%
This is reminiscent of replicator dynamics~\citep{sorin2020replicator}: coordinates with above-average ``fitness'' (larger $\uui{i}$) are amplified relative to those below average, polarizing the mass towards the best-performing coordinates. Motivated by this analogy, we next analyze how softmax normalization induces such polarization and sparsity in $\sos$.

%Here we are leaving some constants and normalization terms for brevity. This differential equality holds similarity to replicator dynamics in population dynamics and evolutionary game theory and dynamics of consesus and polarization of opinions in social networks.  To make the connection clearer, we will take a closer look at the term $\uu  -  \scal{\uu}{\sos} \one$, since $s = \sf(\aa)$ and sums to 1,  $\scal{\uu}{\sos}$ is a weighted mean and the term $\uu  -  \scal{\uu}{\sos} \one$ essentially subtracts the mean from all the co-ordinates. This is analogues to the replicator dynamics where the growth rate of a species is proportional to the difference of its fitness to the average fitness of all species and eventually the fittest survives. It can also be seen as polarization where the opinions drift from their aggregate. Leveraging intutions from the above analogies, we will proceed to study the polarization induced by the softmax normalization.

\paragraph{Repulsion between the coordinates.}
%
To formalize the polarizing effect, we first show that the dynamics induce repulsion between the coordinates of the attention scores $\sos$ and the value projections $\uu$. We begin by stating an assumption on the initialization. 
%
\begin{assumption}\label{ass:init}
  The parameters $\VV, \aa$ are initialized such that $\aa(0) = 0$ and  $\uui{0}(0) > \uui{1}(0) > \ldots > \uui{p-1}(0)$ for the projection $ \uu = \VV^{\top} \bbeta_*$.
\end{assumption}
%
This assumption is mild. The condition $\aa(0)=0$ implies $\sos(0)$ is uniform, which is consistent with standard practice. Without loss of generality, we assume the coordinates of $\uu(0)$ are ordered by re-indexing. The strict ordering holds almost surely under generic random initialization and can be relaxed if needed. We now show that gradient flow preserves this ordering and induces repulsion between coordinates.
%
\begin{restatable}{theorem}{NoCrossingLogistic} \label{thm:no-crossing-logitstic}
 Consider gradient flow on the loss \eqref{eq:logistic-loss-V-a} under the initialization in \cref{ass:init}. Then, for all $t > 0$,
  \begin{enumerate}[label=(\alph*)]
    \item \textbf{Order preservation.} The order of both the projections $\uui{i}(t)$ and attention scores $\sosi{i}(t)$ is preserved
  \begin{align*}
    \uui{0}(t) &> \uui{1}(t) > \ldots > \uui{p-1}(t), \\
    \sosi{0}(t) &> \sosi{1}(t) > \ldots > \sosi{p-1}(t).
  \end{align*}
  \item \textbf{Repulsion.} Pairwise gaps between value projections grow over time: for any $i \neq j$ and  any  $t_+ > t \geq 0$,
  \begin{align*}
     \left( \uui{i}(t_+) - \uui{j}(t_+) \right) > \left( \uui{i}(t) - \uui{j}(t) \right).
  \end{align*}
  \end{enumerate}
\end{restatable}
%
\Cref{thm:no-crossing-logitstic} first establishes a \emph{no-crossing} property: the coordinate ordering is preserved along the trajectory. To prove this, we introduce the Lyapunov potential $\Phi_{ij} = - \left( e^{-\aai{i}} - e^{-\aai{j}}  \right) \left(  \uui{i} - \uui{j} \right)$ for $i\neq j \in [p]$. The derivative of potential is positive, $\der{\Phi_{ij}} > 0 $, hence the relative ordering between the scores and projections is preserved. In particular, $\Phi_{ij}(t) > 0$, for all $t > 0$, hence the co-ordinates never cross and the initial order remains conserved.
%
Second, the gradient flow exhibits repulsion among the coordinates of $\uu$: for $i < j$, we have $\der{ (\uui{i} - \uui{j}) } > 0$, driven by the preserved ordering of the attention scores. 
%
Together, these properties highlight the polarizing nature of the gradient field in the value–softmax model.

\begin{figure*}[t]
    \centering
    \includegraphics[width=\linewidth]{figs/softmax_and_u_evolution_icml_2col.pdf}
    \caption{Experimental verification of Theorem \ref{thm:sparsity-logistic} using logistic loss. The plot shows the evolution of attention scores $\sigma(a)$ and value projections $u$ over time. As predicted, the attention scores converge to a one-hot vector (blue line goes to 1, others to 0) and the value projections $u$ diverge, preserving their order.}
    \label{fig:logistic_evolution}
\end{figure*}


\paragraph{Sparsity of the attention scores.}
%
Having established the polarizing nature of the dynamics, we now quantify its strength and characterize the resulting limiting behavior.
%
\begin{restatable}{theorem}{SparsityLogistic}
\label{thm:sparsity-logistic}
%
Consider gradient flow for the loss \eqref{eq:logistic-loss-V-a} under the initialization in \cref{ass:init}. Let $ \delta = \min_{i < j} \left( \uui{i}(0) - \uui{j}(0) \right) $.  Then we have 
\begin{enumerate}[label=(\alph*)]
    \item \textbf{Vanishing loss.}  
    $ \lim_{t \to \infty} L(\bbeta(t)) = 0 $.
    \item \textbf{Decay of non-maximal scores. } For any  $j \neq 0$ and any $t > 0$,
    \[ \frac{\sosi{j}(t)}{\sosi{0}(t)} \leq \frac{1}{1 + \nicefrac{\delta}{p} \int\limits_{0}^{t} \gamma(\bbeta(s)) ds }. \] 
    \item \textbf{One-hot limit.}
    The polarization coefficient $ \int_{0}^{t} \gamma(\bbeta(s)) ds = \Theta(\log t)$, and consequently the attention scores converge to a one-hot vector, i.e,
    \[ \lim_{t \to \infty} \sosi{0}(t) = 1, \quad \text{and} \quad \lim_{t \to \infty} \sosi{j}(t) = 0, \quad \text{for } j \neq 0. \]
\end{enumerate}
\end{restatable}
%
To prove the theorem, we track the ratio $\sosi{j}(t)/\sosi{0}(t)$ and show that it decays at a rate controlled by the integral $\int_{0}^{t}\gamma(\bbeta(s))\,ds$. 
%
We then prove that this integral diverges as $t\to\infty$, implying that $\sos(t)$ converges to a one-hot vector.
%
The softmax parameterization implicitly biases the dynamics toward low-entropy (peaked) attention.
%
Finally, note that,
\begin{align*}
  \bbeta(t) = \sum_j \sosi{j}(t) \vvi{j}(t).
\end{align*}
Thus, in the limit, $\bbeta(t)$ becomes supported on a single column of $\VV(t)$: among the many possible decompositions of $\bbeta$ as a convex combination of $\{\vvi{j}\}$, gradient flow selects an extremal (one-hot) representation. This provides a theoretical explanation for the empirically observed tendency of softmax attention to become sparse.
%To prove the theorem, we track the ratio of the attention score to the maximum attention score $\sosi{0}(t)$ and show that the ratio decays inversely proportional to the integral of $\gamma(\bbeta(s))$.Next, we show that the integral diverges to infinity as time tends to infinity, which shows that the attention scores converge to a one-hot vector. This shows strikingly that the softmax parameterization --- it implicitly biases the dynamics such that the limiting softmax is one-hot. Note that \begin{align*}  \bbeta(t) = \sum_j \sosi{j}(t) \vvi{j}(t)\end{align*}
%The final embeddings now are only supported on just one column and choose this representations overs all possible ways of representing beta. This provably shows that the softmax parameterizations favors attentions scores with lower entropy and lays theoritical foundations for extensive empirical observations.

\paragraph{Comparison with the dynamics of linear networks.}
%
When $\sigma$ is the identity, the model reduces to a linear network, whose gradient-flow dynamics have been extensively studied in connection with implicit regularization of gradient-based methods. In particular, it is known that the solution converges to a rank-one factorization of $\VV $(see, e.g., \cite{ji2018gradient}). 
%
In our setting, a similar phenomenon occurs: under the softmax parameterization, $\VV(t)$ also becomes approximately rank one (see Appendix~\ref{lem:log-non-maximal-rates}). The crucial difference, however, is that the attention scores $\sos(t)$ converge to a one-hot vector. This behavior does not arise in linear networks, and it forces the dominant right singular direction of $\VV$ to align with a coordinate axis (i.e., to be one-hot) rather than a generic direction.
%

The polarization we observe is not caused by the softmax nonlinearity in isolation, but by the \emph{joint} training of the scores $\aa$ and the value matrix $\VV$. Their co-adaptation drives the dynamics toward low-entropy (peaked) attention patterns.

Finally, it is important not to conflate this effect with sparse regression: the sparsity here is in the attention scores $\sos$, not in the entries of the value matrix $\VV$ or final predictor $\bbeta$.


\subsection{Extensions of the result}
\label{sec:extensions}
%
We now explore how far our analysis extends beyond the logistic-loss setting. In particular, we vary the loss and the attention map to identify when polarization persists and when it fails.
%This section explores how robust our conclusions are to modeling choices made earlier in the paper. We move beyond the logistic-loss setting by varying the loss function and the normalization/nonlinearity in the attention map, highlighting both regimes where the dynamics remain qualitatively similar and where they fundamentally change.
%We now discuss several extensions that illustrate which aspects of our analysis are specific to the logistic loss and which reflect a more general mechanism in the gradient-flow dynamics. In particular, we consider alternative objectives and parametrizations, and examine when the same polarization/sparsification behavior persists or breaks down.
%
\paragraph{Regression.}
We now consider a regression setting in which the target vector $\bbeta_*$ is to be fitted. The loss is 
\begin{align} \label{eq:regression-loss-V-a}
  \LL{\VV, \aa} = \frac{1}{2} \nor{ \bbeta_* - \VV \sf(\aa) }^2.
\end{align}
The associated gradient flow dynamics are
\begin{align*}
  \der{\VV} &=  \left(\bbeta_* - \VV \sos \right) \sos^{\top} \der{t}, \\
  \der{\aa} &=  \left[ \diag{\sos} - \sos \sos^{\top} \right] \VV^{\top} \left(  \bbeta_* - \VV \sos \right) \der{t}.
\end{align*}
We make the following initialization assumption.
\begin{assumption} \label{assm:regression-init}
  The parameters $\VV, \aa$ are initialized such that $\VV(0) = 0$ and the attention scores satisfy $ \sosi{0}(0) > \sosi{1}(0) > \ldots > \sosi{p-1}(0) $.
\end{assumption}
%
Initializing the linear predictor to zero is standard in analyses of gradient flow for linear regression and simplifies the resulting dynamics. We assume the attention scores are strictly ordered at initialization to avoid ties and break symmetry. Without loss of generality, we label the coordinates so that the ordering is decreasing. Under this initialization, $\VV(t)$ remains rank one, which allows us to reduce the dynamics to a coupled system similar to \cref{eq:logistic-loss-u-a}, as stated next.
%it takes the form $\VV = \frac{\bbeta_*}{\|\bbeta_*\|^2} \uu^{\top} $, using this we can again cast the dynamics similar to Eq.~\eqref{eq:logistic-loss-u-a}, however with $$ \gamma(\bbeta) =  1 - \frac{\scal{\uu}{\sos}}{\nor{\bbeta_*}^2}, $$defined accordingly. 
\begin{restatable}{lemma}{Regression}
Consider gradient flow for the loss \eqref{eq:regression-loss-V-a} under \cref{assm:regression-init}. Then:
\begin{enumerate}[label=(\alph*)]
    \item \textbf{Rank-one structure.} 
 There exists $\uu(t)\in\R^p$ such that $\VV = \frac{\bbeta_*}{\|\bbeta_*\|^2} \uu^{\top} $.
    \item \textbf{Reduced dynamics.} The pair $\aa, \uu $ satisfies  
\begin{align} \label{eq:regression-loss-u-a}
  \mdot{\uu} &= - \gamma(\bbeta) \, \nor{\bbeta_*}^2 \sos, \\
  \mdot{\aa} &= - \gamma(\bbeta) \, \left[ \diag{\sos} - \sos \sos^{\top} \right] \uu, 
\end{align}
where $ \gamma(\bbeta) =  1 - \frac{\scal{\uu}{\sos}}{\nor{\bbeta_*}^2}$.
\end{enumerate}
\end{restatable}
%
Using the same techniques as in the logistic-loss case, we can derive an analogue of \cref{thm:no-crossing-logitstic} for the regression dynamics.
%
A key difference is that the cumulative polarization strength is finite:  $\int_{0}^{\infty} \gamma(\bbeta(s)) \der{s} = \Theta(1)$ (see Appendix for the result).
Consequently, polarization is only partial and $\sos(t)$ does not converge to a one-hot vector in general. More broadly, sparsification is controlled by the integral of the  gradient magnitude: the slower the gradient decays, the stronger the induced sparsity in the attention scores.
%
This suggests that ill-conditioned regression problems (e.g., $\| \bbeta_* - X \VV \sf(\aa) \|^2$ for a matrix  $X$ ) may exhibit stronger sparsification, since ill-conditioning slows down convergence. We empirically verify this relationship in Figures \ref{fig:main-cond-numbers} and \ref{fig:condition_numbers}.
%\todo{Note however that adding noise in the dynamics as in \cite{vivien2022label} does not solve the problem since tbd.} 
% On these lines, one can consider ill-conditioned regression problems for example $\| \bbeta_* - X \VV \sf(\aa) \|^2$, for some matrix $X$ and once can see the co-relation between the condition number sparsification of the attention scores and the condition number as for ill-conditioned problems the speed of gradient convergence is slower and hence the solutions are sparser, see Figure.~\ref{}. 

% \subsection{Some generalizations.}\paragraph{Other loss functions.} Other loss functions of choice could be KL-divergence, where $\ell(\bbeta) = - \scal{p_*}{\log(\bbeta)}$ where $p_*$ is the target probability vector. The gradient flow dynamics is similar to the case of regression and also polarization of depleting strength can be observed. 

\begin{figure}[b]
  \centering
  \includegraphics[width=\columnwidth]{figs/condition_numbers_small.pdf}
  \caption{Evolution of attention logits in a regression problem with condition number $\kappa = 1$ and 5.}
  \label{fig:main-cond-numbers}
\end{figure}


\paragraph{ KL divergence.}
A natural alternative is the KL divergence between a target distribution $p_*$ and a predicted distribution $\bbeta$, e.g., $\ell(\bbeta) = -\langle p_*, \log \bbeta\rangle$, where the logarithm is applied elementwise. The resulting gradient-flow dynamics are analogous to the regression case, and we again observe polarization, albeit with diminishing strength compared to the case of classification (Figure \ref{fig:kl}).

\paragraph{Other non-linearities and normalization functions.} 
%
We first consider replacing softmax with elementwise nonlinearities such as the sigmoid activation  $\sigma_{sig}(a)_i = \nicefrac{1}{1 + \exp(-a_i)}$ or the Relu. For both the regression and logistic losses, we do not observe the same polarization behavior under sigmoid (and similarly under ReLU), see Figure~\ref{fig:toy-many-plots} for reference. Intuitively, the key mechanism behind polarization in our analysis is the mean-centering term that appears in the softmax Jacobian; elementwise nonlinearities such as sigmoid or ReLU do not induce an analogous ``fitness minus average fitness'' interaction across coordinates.

We next consider alternative \emph{normalization} maps, which retain this interaction structure. Let $f$ be a monotonically increasing function and define the normalization $\sf_{f}: \R^{p} \to \R^{p}$ by
\begin{align*}
  \sf_{f}(a)_i = \frac{f(a_i)}{\sum_{j=1}^{p} f(a_j)}.
\end{align*}
%
Writing $\sos_f \coloneqq \sf_f(\aa)$, the gradient flow dynamics for the logistic loss take the form
\begin{align*}
  \der{\uu} &=  \gamma(\bbeta) \, \|\bbeta_*\|^2   \sos_{f} \, \der{t}, \\
  \der{\aa} &=  \gamma(\bbeta) \, \frac{f'(\aa)}{\sum_{j=1}^{p} f(\aa_j)} \odot ( \uu - \scal{\uu}{\sos_f} \one ) \, \der{t},
\end{align*}
%
This retains the replicator-type structure, and an analogue of \cref{thm:no-crossing-logitstic} can be established. 
%
The polarization strength now depends on $f$ and $f'$. For instance, with $f(x)=x$ (which does not enforce positivity), the polarization is typically insufficient to yield one-hot scores, where as with $f(x)=x^2$ the dynamics can drive $\sos_f$ toward a one-hot limit, see Figure~\ref{fig:toy-many-plots} provides empirical evidence for these behaviors.
%

\paragraph{Limitations.}
%
A key limitation of our analysis is that, for the losses considered above, the gradient $\nabla \ell(\bbeta)$ is aligned with the target direction $\bbeta_*$ (in the regression case, this relies on \cref{assm:regression-init}). 
%
This alignment substantially simplifies the dynamics and enables the reduction to the two-vector system in $\aa$ and $\uu$. 
%
In practice, such an alignment need not hold, and the present analysis may therefore not apply directly. Nevertheless, our experiments indicate that the polarization effect persists even without exact alignment; see \cref{fig:main-cond-numbers} as alignment breaks in the large condition number case.
 
\section{Implications of softmax-induced polarization for attention sparsity and token influence}
\label{sec:implications}
We now connect our theoretical polarization result to qualitative behaviors observed in attention mechanisms. In particular, we discuss how softmax-induced sparsity can lead to attention sinks (and associated large activations) and to an imbalance in how individual tokens influence the model’s predictions. The code is available at \url{https://github.com/tml-epfl/softmax}



\subsection{Attention sinks and massive activations}
\paragraph{Attention sinks.} 
Attention sinks~\citep{xiao2024efficient} are attention patterns in which the maximal attention is over a specific token (often special tokens such as BOS or delimiter/punctuation tokens).
%
Attention sinks are suggested to be formed when the heads need to be turned off or when they have to create a explicit bias due to absence of biases in self-attention architecture~\citep{sun2024massive,zhang2025attention,gu2025when}. 
%
However, to create a bias, there is no reason to focus all attention on a single token; a distributed or uniform pattern over a set of tokens could also create such a bias.
%
%To give some insights into this mechanism, consider the value-softmax model and the results from Theorem~\ref{thm:sparsity-logistic}.  Lets say the output of the attention block is $\sf(\bA)\VV$, where $\bA \in \R^{T \times T}$ is the attention scores matrix with $\sf$ applied row-wise.
%
%\begin{align*}
%  \bbeta &= \sf(\bA) \VV, \\
%  \LL{\VV, \bA} = \ell(\bbeta) &= \frac{1}{T} \sum_{i=1}^{T} \log\left( 1 + \exp\scal{\bbeta_*}{- \bbeta[t]} \right).
%\end{align*}
%Here, the objective is to create the same bias $\bbeta_*$ for all the tokens in the sequence. Using Theorem~\ref{thm:sparsity-logistic}, we can see that the attention scores for each token converges to a \textbf{same} one-hot vector determined by the random initialization of $\VV$. Hence for each token the attention is dumped on a single token, which creates an attention sink. 
%This provides some intution that the attention sinks \textit{could} be a consequence of the optimization dynamics of the softmax parameterization. 

\begin{table}[t]
\centering
\small
\setlength{\tabcolsep}{6pt}
\renewcommand{\arraystretch}{1.35}

\resizebox{\linewidth}{!}{%
\begin{tabular}{l | l | l}
\toprule
Name & $\mathrm{sim}(\phi(\mathbf q_i), \phi(\mathbf k_j))$ & Normalization \\
\midrule
softmax &
$\displaystyle \exp\!\left(\frac{\mathbf q_i \mathbf k_j^\top}{\sqrt{d_h}}\right)$ &
$\displaystyle \sum_{j'=1}^{m}\exp\!\left(\frac{\mathbf q_i \mathbf k_{j'}^\top}{\sqrt{d_h}}\right)$ \\

sigmoid &
$\displaystyle \sigma\!\left(\frac{\mathbf q_i \mathbf k_j^\top}{\sqrt{d_h}}\right)$ &
$\displaystyle \sum_{j'=1}^{m}\sigma\!\left(\frac{\mathbf q_i \mathbf k_{j'}^\top}{\sqrt{d_h}}\right)$ \\

elu &
$\displaystyle \mathrm{elu}\!\left(\frac{\mathbf q_i \mathbf k_j^\top}{\sqrt{d_h}}\right)+1$ &
$\displaystyle \sum_{j'=1}^{m}\Big(\mathrm{elu}\!\left(\frac{\mathbf q_i \mathbf k_{j'}^\top}{\sqrt{d_h}}\right)+1\Big)$ \\

elu $\cdot$ elu &
$\displaystyle \frac{\big(\mathrm{elu}(\mathbf q_i)+1\big)\big(\mathrm{elu}(\mathbf k_j)+1\big)^\top}{\sqrt{d_h}}$ &
$\displaystyle \sum_{j'=1}^{m}\frac{\big(\mathrm{elu}(\mathbf q_i)+1\big)\big(\mathrm{elu}(\mathbf k_{j'})+1\big)^\top}{\sqrt{d_h}}$ \\

mlp &
$\displaystyle \frac{\mathrm{mlp}(\mathbf q_i)\,\mathrm{mlp}(\mathbf k_j)^\top}{\sqrt{d_h}}$ &
$\displaystyle \max\!\left(\left|\sum_{j'=1}^{m}\frac{\mathrm{mlp}(\mathbf q_i)\,\mathrm{mlp}(\mathbf k_{j'})^\top}{\sqrt{d_h}}\right|,\,1\right)$ \\

linear &
$\displaystyle \frac{\mathbf q_i \mathbf k_j^\top}{\sqrt{d_h}}$ &
$\displaystyle \max\!\left(\left|\sum_{j'=1}^{m}\frac{\mathbf q_i \mathbf k_{j'}^\top}{\sqrt{d_h}}\right|,\,1\right)$ \\

\bottomrule
\end{tabular}%
}
\vspace{2pt}
\caption{Attention functions we test in Section \ref{sec:implications}. When computing the normalization, $m =$ sequence length for non-causal models and $m = i$ for causal models. For all functions but softmax, we test their variants both with and without normalization.}
\label{table:attention-types}
\end{table}

To shed light on why optimization may nevertheless favor sinks, consider the value--softmax model in which the attention output is
$\bbeta = \sf(\bA)\VV$, where $\bA\in\mathbb{R}^{T\times T}$ is the attention scores matrix (softmax applied row-wise) and $\VV\in\mathbb{R}^{T\times d}$ are value vectors.
Suppose we optimize $(\VV,\bA)$ so that each position $t$ produces the same target vector $\bbeta_*\in\mathbb{R}^d$, e.g.,
\begin{align*}
  \LL{\VV, \bA} = \ell(\bbeta) = \frac{1}{T} \sum_{i=1}^{T} \log\left( 1 + \exp\scal{\bbeta_*}{- \bbeta[t]} \right).
\end{align*}
Extending \cref{thm:sparsity-logistic}, the row-wise softmax weights converge to a one-hot vector supported on the initially maximal coordinate of $\uu(0)=\VV^\top\bbeta_*$.
Under Assumption~\ref{ass:init} this coordinate is $0$ for all rows, so $\sf(\bA[t,:])\to e_0$ for every $t$, yielding an attention sink at token $0$. This result suggests that attention sinks \emph{can} emerge as a consequence of the optimization dynamics induced by the softmax parameterization.
%
%

\paragraph{Massive activations.} 
%
\citet{sun2024massive} has reported that attention sinks are also coupled with massive activations in some feature dimension, i.e., a regime where a small fraction of activations takes values that are significantly larger than the rest.
%
A reformulation of the value--softmax viewpoint provides intuition for how this can arise from the same polarizing dynamics.
% 
Recall that the forward pass for some intermediate residual state $R \in T \times T$, and the residual state of the last token $r$ is given by 
\begin{align*}
  R_{+} &=   \sf\left( R  QK^{\top} R^{\top} \right) R V, \\
  r_{+} &=  V^{\top} R^{\top} \sf( RKQ^{\top}z). 
\end{align*}
%Note that here the product is between the residual state matrix and the softmax outputs, similar to stylized value-softmax model. One can consider the simple model such as 
Note that here the product is between the residual state matrix and the softmax outputs, similar to the value-softmax model with an exception that the softmax inputs are now not free. One can express this idea with a formulation
$\LL{\RR, \aa}= \bloss\left( \RR \sf(\RR \aa) \right)$, where $\bloss$ is a logistic loss function, and see that this model qualitatively follows the behavior of the value-softmax model. 
%\todo{not the norm of U but norm of one column of u vs others} 
In Figure \ref{fig:massive_activations}, we can see the polarization in softmax scores and also the polarization in the embeddings of the residual state for this model. This shows that the attention sinks and massive activations  can arise from underlying the polarizing dynamics. 

\paragraph{Induction heads acting as attention sinks.}

\begin{figure}[t]
  \centering
  \includegraphics[width=\columnwidth]{figs/sink_proportions_aux0.2_shift_add.pdf}
  \caption{Proportion of sink heads at the 2nd layer of Transformer trained for the induction task.}
  \label{fig:main-exp-induction}
\end{figure}

To test whether the results proven for the value-softmax model hold for more general Transformer setups, we run additional experiments in more realistic transformer model (details in Appendix \ref{app:experimental-details}). 
First, we create a simple setup where attention sinks emerge. We train 2-layer Transformers with a variable number of attention heads on an induction task: predict a bi-gram learned in-context. Our task is a simple variant of Bigram-Backcopy task of  ~\citep{guo2024active,bietti2024birth}. Each input object is a sequence of tokens sampled uniformly at random, 
%where no informed prediction is possible, 
followed by a bigram, that can be predicted by a construction known as \emph{induction heads} \citep{olsson2022context}. Example input is shown below:
\begin{figure}[h]
\vspace{-4pt}
\centering
{
\textcolor{black!70}{\textbf{BOS}}\;
\textcolor{black!70}{4}\;
\textcolor{black!70}{2}\;
\textcolor{black!70}{6}\;
\textcolor{blue!70!black}{12}\;
\textcolor{green!60!black}{8}\;
\textcolor{black!70}{11}\;
\textcolor{blue!70!black}{12$^\ast$}\;
\textcolor{green!60!black}{8}
}
\vspace{-15pt}
\end{figure}

At the positions where tokens are sampled at random,  an induction heads need not appear and they can act as a fixed attention bias similar to the 
%
%The induction heads are not useful at the positions of random tokens, which %incentivizes models to suppress them, forming the conditions for
the emergence of attention sinks as described in Section \ref{sec:implications}. 
On this task, we evaluate the formation of attention sink and its dependence on various choices of functions listed in Table~\ref{table:attention-types} as an alternative to compute attention. 
We define an attention head as a \emph{sink head} if, at positions where tokens are sampled uniformly at random, it assigns more than 90\% of the total attention weight to the first token. We report the average proportion of such heads in Figure \ref{fig:main-exp-induction}.
% We consider an attention head a \emph{sink head} if it assigns more than 90\% of total attention weight at position where tokens are sampled uniform at random on the first token, and report the average proportion of such heads in Figure \ref{fig:main-exp-induction}. 
% Following \citet{gu2025attention}, we consider a range of attention functions alternative to softmax, listed in Table \ref{table:attention-types}.
In line with the results in Section \ref{sec:extensions}, the normalized positive attention variants are associated with consistently higher proportion of sink heads than their unnormalized counterparts or non-positivite variants (mlp, linear).% \todo{@Aditya does it make sense?}

\begin{figure}[t]
  \centering
  \includegraphics[width=\columnwidth]{figs/classification_with_sample_tracking.pdf}
  \caption{Evolution of the attention scores for two differently labeled samples in the classification experiment.}
  \label{fig:exp-classifcation-attentions}
\end{figure}

\begin{figure}[b]
  \centering
  \includegraphics[width=\columnwidth]{figs/fliprate_vocab41_mlp_ln.pdf}
  \caption{Success rates of adversarial token 1-token flips in 1-layer 1-head Transformer trained for classification.}
  \label{fig:main-exp-classification}
\end{figure}


%\paragraph{Difference in attention sparsity between pretrained softmax and sigmoid LLMs. }
% Additionally, we check if the observations achieved in toy setups hold for real-world Large Language Models. We compare softmax and sigmoid (unnorm.) 7-billion-parameter LLMs matched by hyperparameters and compute, released by \citet{ramapuramtheory}. We run a forward pass of those models on random samples from the Pile dataset \cite{gao2020pile} and measure the sparsity score of each head as the average proportion of total attention weight assigned to the max-logit token. Figure \ref{fig:main-llm-pretrained} shows the distribution of this metric for all heads in both models, supporting the earlier conclusion that softmax leads to much sparser attention outputs and higher chance of sink formation.

\paragraph{Pretrained LLMs.}
Additionally, we verify whether the observations from toy setups hold for real-world large language models. We compare softmax and sigmoid (unnormalized) 7-billion-parameter LLMs matched by hyperparameters and compute, released by \citet{ramapuramtheory}. We run a forward pass of both models on random samples from the Pile dataset \cite{gao2020pile} and measure the sparsity score of each head as the average proportion of total attention weight assigned to the max-logit token. Figure \ref{fig:main-llm-pretrained} shows the distribution of this metric across all heads in both models, confirming that softmax leads to significantly sparser attention outputs and a higher likelihood of sink formation. This result aligns with \citet{gu2025when}, who reported similar behaviors in models with 1B parameters.
%This observation corroborating the finding of  \citet{gu2025when} on LLMs with 1B parameters. 

% To verify if our observations from stylized models generalize to realistic settings, we analyzed 7-billion-parameter LLMs using both softmax and sigmoid attention (without normalization).  (matched for compute and hyperparameters from \citet{ramapuramtheory}). We performed forward passes on random samples from the Pile dataset \cite{gao2020pile} to calculate a 'sparsity score' for each head, defined as the average attention weight allocated to the max-logit token. The resulting distributions, shown in Figure \ref{fig:main-llm-pretrained}, corroborate our earlier findings: softmax attention produces significantly sparser outputs and a higher propensity for sink formation compared to sigmoid.

%\paragraph{Imbalance in token influence.}
\subsection{Imbalance in token influence}

% The second major implication of our result is the emergence of tokens that are anomalously important for the model's output. If, as dictated by Theorem~\ref{thm:sparsity-logistic}, attention weights of a certain attention head converge to a one-hot vector, then all but one token do not influence the output at all, since it depends on the max-logit token only. Thus, the prediction becomes particularly sensitive to the perturbations of that one token. We verify this below.

A second major implication of our results is the emergence of tokens that exert a disproportionate influence on the model's output. As predicted by Theorem~\ref{thm:sparsity-logistic}, when attention weights converge to a one-hot vector, the output effectively ignores all context except for the single 'max-logit' token. Consequently, the model's predictions become highly sensitive to perturbations of that specific token. % We empirically verify this phenomenon below.

%motivation for the classification experiment here.%{\color{red} should it be a separate paragraph in implications or just a part of the experiment description?}

%\subsection{Experiments beyond the value-softmax model}
%\label{sec:experiments-beyond}

%To test whether the results proven for the value-softmax model hold for more general Transformer setups, we additionally run experiments for more realistic settings\footnote{Additional details on the experimental settings and results are provided in Appendix TODO.}.

%\paragraph{Emerging imbalance in token influence on classifier's predictions.}

We construct a simple classification dataset in the following way:
Consider a vocabulary $V$ and partition $V$ into $k$ sets $V^{(1)},\ldots,V^{(K)}$ of equal size. An input and its corresponding label are  
$x_i = v^{(k)}_{(i,1)},\dots,v^{(k)}_{(i,n)}, y_{i}=k$, where $x_i$ is a sequence of $n$ tokens sampled from $V^{(K)}$ and labeled $k$. Note that as the sets are disjoint any one of the input tokens unambiguously predicts the label, thus any form of attention map allows to solve the classification problem.

Due to the low-entropy bias of softmax, however, we expect that the attention weights will concentrate on few tokens, effectively ignoring the others. We test this by examining the behavior of the classifier in adversarial out-of-distribution setting where we are allowed to change one token in the input with the goal of flipping the classifier's prediction. The success rates of this experiment are presented in Figure \ref{fig:main-exp-classification}. As expected, the softmax model's predictions can be easily reversed, however the effect gradually vanishes as more than 1 head and layer are added to the model, as different heads, although sparse, concentrate attention on distinct tokens.

%We verify the token influence imbalance described above. 
% We construct a simple classification dataset by evenly splitting the vocabulary $V$ into $K$ disjoint sets $V^{(1)},\dots,V^{(K)}$ and sampling pairs $x_i = v^{(k_i)}_{j_{i1}},\dots,v^{(k_i)}_{j_{in}}, y_i=k$. Any one of the input tokens unambiguously predicts the label, thus any form of attention map allows to solve the classification problem.

% Due to the low-entropy bias of softmax, however, we expect that the attention weights will concentrate on few tokens, effectively ignoring the others. We test this by examining the behavior of the classifier in adversarial out-of-distribution setting where we are allowed to change one token in the input with the goal of flipping the classifier's prediction. The success rates of this experiment are presented in Figure \ref{fig:main-exp-classification}. As expected, the softmax model's predictions can be easily reversed, however the effect gradually vanishes as more than 1 head and layer are added to the model, as different heads, although sparse, concentrate attention on distinct tokens.

\begin{figure}[t]
  \centering
  \includegraphics[width=\columnwidth]{figs/sink_scores.pdf}
  \caption{Distribution of attention head sparsity scores for pretrained softmax/sigmoid LLMs.}
  \label{fig:main-llm-pretrained}
  \vspace{-2pt}
\end{figure}

\section{Conclusion}
%We have provably shown the softmax normalization leads to solutions where the softmax output has low entropy. We have used some principled simplifications of the attention block and some typical loss functions to prove the result. We have further validated our observations on the simple model with more realistic architechtures.   \todo{empirically validated with experiments}.
%We show that the softmax normalization in attention induces an implicit optimization bias toward low-entropy (peaked) weight vectors when the scores and values are trained jointly under gradient flow. Using a simplified attention model and standard losses, we prove polarization dynamics that drive the softmax output toward extremal solutions, and we characterize how this effect depends on objective and convergence behavior. Finally, we corroborate the theory empirically in more realistic architectures, demonstrating that normalization choice modulates attention sparsity and sink-like behavior.


This work isolates an implicit optimization bias in the value–softmax parameterization $\VV\sf(\aa)$: under gradient flow, jointly training the scores and values  polarizes the attention distribution toward low-entropy solutions, even when many higher-entropy decompositions represent the same predictor. In a stylized classification setting, we prove that $\sf(\aa)$ converges to a one-hot vector, selecting an extremal representation. We extend the analysis to regression and related objectives, clarifying how convergence speed and activation choices modulate polarization. 
Finally, we connect this dynamics to transformer behavior and identify its connection with attention sinks and associated pathologies. %We support the theory with experiments showing how normalization impacts sink formation in induction-trained models and how pretrained softmax LLMs exhibit  higher attention sparsity than sigmoid variants.
We support the theory with experiments showing how attention type impacts sparsity in models trained for induction and classification, as well as for pretrained LLMs.

\bibliographystyle{abbrvnat}
\bibliography{references}

\newpage
\appendix
\onecolumn




\section{Some Properties of the Value-softmax Model}

\paragraph{Notation.} For any function $f : \R^{m} \to  \R^{n}$, for $x \in \R^{m}$, denotes $\jacb{x}{f} \in \R^{n \times m}$ denotes Jacobian of the function. 


\paragraph{Gradient Flow Dynamics.} We will consider the gradient flow dynamics of 
the loss function $\loss(\VV, \aa)$ given by
\begin{align} \label{eq:gradient-flow}
    \frac{d\VV}{dt} &= - \nabla_{\VV} \loss(\VV, \aa),  \quad 
    \frac{d\aa}{dt} = - \nabla_{\aa} \loss(\VV, \aa).
\end{align}
We will use the chain rule to compute the gradients of $\loss(\VV, \aa)$ as follows:
\begin{align*}
\jacb{\VV}{\loss(\VV, \aa)} &= \jacb{\bbeta}{\ell} \jacb{\VV}{\bbeta(\VV, \aa)}, \\ 
\jacb{\aa}{\loss(\VV, \aa)} &= \jacb{\bbeta}{\ell} \jacb{\aa}{\bbeta(\VV, \aa)}.
\end{align*}
Also, with an abuse of notation, we also denote the jacobian of $\bbeta$ with respect to $\sos$.  
\begin{align*}
    \jacb{\sos}{\loss(\VV, \aa)} = \jacb{\bbeta}{\loss} \jacb{\sos}{\bbeta(\VV, \aa)}, \\
\end{align*}
% We can compute the Jacobians as follows:
% \begin{align}
Using the properties of Jacobians in Property~\ref{prop:jacobians} and using the notation, 
\begin{align*}
 \rrr^{\top} = - \jacb{\bbeta}{\loss}, \text{ i.e., } \quad \rrr = - \nabla_{\bbeta} \loss. 
\end{align*}
We can write the Jacobians as follows, 
\begin{align*}
    \jacb{\VV}{\loss(\VV, \aa)} &=  \jacb{\bbeta}{\ell} \, \, \left[ \sos^{\top} \otimes \I_{d} \right] = - \rrr^{\top}  \left[ \sos^{\top} \otimes \I_{d} \right] = - \sos^{\top} \otimes \rrr^{\top}.   \\
    \jacb{\aa}{\loss(\VV, \aa)} &=  \; \jacb{\bbeta}{\ell} \, \left[ \VV \left( \diag{\sos} - \sos \sos^{\top} \right) \right] = - \rrr^{\top} \left[ \VV \left( \diag{\sos} - \sos \sos^{\top} \right) \right]. \\
\end{align*} 
With $\rrr = - \nabla_{\bbeta} \loss$,  the jacobians now can be written as gradients in the following way:
\begin{align} \label{eq:gradients-general-loss}
    \nabla_{\VV} \loss(\VV, \aa) &=  - \rrr \sos^{\top},  \\
    \nabla_{\aa} \loss(\VV, \aa) &=  - \left( \diag{\sos} - \sos \sos^{\top} \right) \VV^{\top} \rrr.
\end{align}
\section{Logistic Regression}

In this section, we will analyze the gradient flow dynamics of the value-softmax model over the logistic loss function. First, let us define the logistic loss function over the parameter $\bbeta \in \R^{p}$ as follows:
\begin{align} \label{eq:logistic_beta}
    \bloss(\bbeta) = \log \left( 1 + \exp\left( -  \scal{\bbeta_*}{\bbeta}  \right) \right).
\end{align}
This just changes the scale of the problem and does not affect the analysis. 
\paragraph{Gradient flow over logistic loss.} The gradient of the logistic loss function is given as follows:
\begin{align*}
    - \nabla_{\bbeta} \bloss(\bbeta) = \frac{1}{1 + \exp\left(  \scal{\bbeta_*}{\bbeta}  \right)} \bbeta_*.
\end{align*}
Note that it is always along the direction of $\bbeta_*$ and the scale depends on the parameter $\bbeta$.  Denote the scalar term in the front of $\bbeta_*$ as $\gamma(\bbeta)$, i.e.,
\begin{align*}
    \gamma(\bbeta) \coloneqq \frac{1}{1 + \exp\left(  \scal{\bbeta_*}{\bbeta}  \right)}.
\end{align*}
From now, we will drop the $\bbeta$ and denote it just by $\gamma$.  Using the gradient computation in Eq.~\eqref{eq:gradients-general-loss}, the gradient flow dynamics of $\VV$ and $\aa$ writes:
\begin{align*}
    \frac{d\VV}{dt} &=  \gamma \, \bbeta_* \sos^{\top},  \quad 
    \frac{d\aa}{dt} =  \gamma \, \left( \diag{\sos} - \sos \sos^{\top} \right) \VV^{\top} \bbeta_*.
\end{align*}
The evolution of $\sos$ can be written as:
\begin{align*}  
        \frac{d\sos}{dt} &=  \gamma \, \left( \diag{\sos} - \sos \sos^{\top} \right)^2 \VV^{\top} \bbeta_*.
\end{align*}
Tracking the evolution of the projection vector $\uu$ given by $\uu \coloneqq \VV^{\top} \bbeta_*$, we have,
\begin{align*}
    \frac{\der{\uu}}{\der t} &= \gamma \sos \nor{\bbeta_*}^2.
\end{align*}
This gives us the system of equations governing the gradient flow dynamics of the value-softmax model over the logistic loss function as follows:
\begin{align} 
  \mdot{\uu} &= - \gamma(\bbeta) \, \nor{\bbeta_*}^2 \sos, \\
  \mdot{\aa} &= - \gamma(\bbeta) \, \left[ \diag{\sos} - \sos \sos^{\top} \right] \uu.
\end{align}

%     \frac{d\vvc{i}}{dt} &=   \gamma \, \sos_i \, \bbeta_*,  
%     \frac{d\aai{i}}{dt} &=  \gamma \left[ \sosi{i} \uui{i} - \gamma \sosi{i} \scal{\sos}{\uu} \right], \\
%     \frac{d\sosi{i}}{dt} &=  \gamma \left[ \sosi{i}^2 \left(  \uui{i} - \scal{\sos}{\uu} \right) - \Gamma(\sos, \uu) \sosi{i} \right]. 
% \end{align*}
\paragraph{No-crossing property.} Now that we have established the gradient flow dynamics of the value-softmax model over the logistic loss function, we will analyze some properties of the dynamics. We will now give a lemma which highlights the "no-crossing" property of the dynamics. The lemma proceeds by studing the differences between co-ordinates of the projection vector $\uu$ and $\aa$.



% \begin{assumption}[Initialization]
%     We choose and initialization such that $\aa(0) = 0$ and with a loss of generality $\VV(0)$
%     such that 
%     $$ \uui{1}(0) = \scal{\vvc{1}(0)}{\bbeta_*} > \ldots> \uui{i}(0) = \scal{\vvc{i}(0)}{\bbeta_*} > \uui{d}(0) = \scal{\vvc{d}(0)}{\bbeta_*} $$.
% \end{assumption}

% \begin{lemma}
% \textbf{No-Crossing.} Under the above initialization assumption, for any $i < j$, we have,
% \begin{align*}
%     \aai{i}(t) > \aai{j}(t), \quad \sosi{i}(t) > \sosi{j}(t), \quad \uui{i}(t) > \uui{j}(t), \quad \forall t \geq 0.
% \end{align*}    
% \end{lemma}

\NoCrossingLogistic*
\begin{proof}
    Let $     \cA = \sum_{i} \exp{\aai{i}} $, and note that $\sosi{i}{=} \, \nicefrac{\exp{\{\aai{i}\}}}{\cA}$. Using this, we can write the evolution of $\aai{i}$ as follows:
\begin{align*}
    \frac{d\aai{i}}{dt} &=  \gamma \frac{\exp\{\aai{i}\}}{\cA}  \left[ \uui{i} - \gamma \scal{\sos}{\uu} \right], \\
    \frac{d\{ {-} \exp\{{-\aai{i}} \}\}}{dt} &=   \frac{\gamma }{\cA}  \left[ \uui{i} - \gamma \scal{\sos}{\uu} \right].
\end{align*}
Taking the differences between the evolution of two co-ordinates $i$ and $j$, we have,
\begin{align*}
    \frac{d \{ {-}\left( \exp{-\aai{i}} - \exp{-\aai{j}} \right) \}}{dt} &=   \frac{\gamma }{\cA}  \left[ \uui{i} - \uui{j} \right], \\
    \frac{ d \{ \uui{i} - \uui{j} \} }{ d t} &=  \gamma \, \left( \frac{\exp{\{\aai{i}\}} - \exp{\{\aai{j}\}}}{\cA} \right) \, \| \bbeta_* \|^2. 
\end{align*}
Notice how the differences reinforce each other, to formulate this we will consider the following Lyapunov function. For $i,j \leq p$,
\begin{align*}
    \Phi_{ij}(t) = (\uui{i}(t) - \uui{j}(t)) \{ {-}\left( e^{-\aai{i}(t)} - e^{-\aai{j}(t)} \right) \}.
\end{align*}

Consider the evolution of the potential $\Phi_{ij}(t)$, note that, due to the  
\begin{align*}
    \frac{d}{dt} \left[ (\uui{i} - \uui{j}) \{ {-}\left( e^{-\aai{i}} - e^{-\aai{j}} \right) \}  \right] &= \frac{\gamma}{\cA} \left( (\uui{i} - \uui{j})^2 \right) + \frac{\gamma}{\cA} \| \bbeta_* \|^2 \{e^{\aai{i}} - e^{\aai{j}}\} \{ {-}\left( e^{-\aai{i}} - e^{-\aai{j}} \right) \},  \\
    \frac{d}{dt} \left[ (\uui{i} - \uui{j}) \{ {-}\left( e^{-\aai{i}} - e^{-\aai{j}} \right) \}  \right] &= \frac{\gamma}{\cA} \left( (\uui{i} - \uui{j})^2 \right)
     + \frac{\gamma}{\cA} \| \bbeta_* \|^2 \frac{ \left(e^{\aai{i}} - e^{\aai{j}} \right)^2 }{e^{\aai{i}} e^{\aai{j}} }.
\end{align*}

    Note that at $t=0$, we have $\aai{i}(0) = \aai{j}(0) = 0$, $\sosi{i}(0) = \sosi{j}(0) = {1}/{p}$ and $\uui{i}(0) > \uui{j}(0)$. So $\Phi_{ij}(0) = 0$ and $\Phi_{ij}(t) > 0$, for any time $t > 0$. The strict inequality comes from the gradient at $t{=}0$ is non-zero. 

    Note that for any $i,j$, we have, $\Phi_{ij} > 0$. Hence, there cannot exist a time $t$ such that $\uui{i}(t) = \uui{j}(t)$ or $\aai{i}(t) = \aai{j}(t)$. Hence the attention scores and the projections cannot cross each other. Hence the order  $\uui{0} > \uui{1} > \uui{2} > \ldots $ at initialization is preserved for all time and also the order $\aai{0} > \aai{1} > \aai{2} > \ldots $ is enforced at time $t > 0$ due to the positivity of $\Phi_{ij}(t)$.
   
% Note that at $t=0$, we have,
% \begin{align*}
%     (\uui{i}(0) - \uui{j}(0)) \{ {-}\left( e^{-\aai{i}(0)} - e^{-\aai{j}(0)} \right) \} = 0.
% \end{align*}
% This implies that for any $t \geq 0$, we have,
% \begin{align*}
%     (\uui{i}(t) - \uui{j}(t)) \{ {-}\left( e^{-\aai{i}(t)} - e^{-\aai{j}(t)} \right) \} \geq 0,
% \end{align*}
% and complete the proof. 
% \end{proof}
The other implication of this no-crossing property is that the differences increase with time. 
For any $i < j$, we have,
\begin{align*}
\frac{ d \{ \uui{i} - \uui{j} \} }{ d t} &=  \gamma \, \left( \frac{\exp{\{\aai{i}\}} - \exp{\{\aai{j}\}}}{\cA} \right) \, \| \bbeta_* \|^2 > 0, \text{ for } t > 0. 
\end{align*}
Hence, $\uui{i} (t_+) - \uui{j} (t_+) > \uui{i} (t) - \uui{j} (t) $ for any $t_+ > t > 0$. A similar argument holds for $\aai{i}$ and $\sosi{i}$.
\end{proof}


\paragraph{Sparsity of the attention.} In the previous lemma, we established the no-crossing property of the dynamics and the repulsive forces. In this section, we will analyze the strength of this repulsive feild to establish the sparsity of the attention scores. To begin with, we will analyze the ration of the scores of two co-ordinates $i$ and $j$. Along with this, we will assume that $$\nor{\bbeta_*} = 1.$$ 


\SparsityLogistic*
\begin{proof}
 To prove this theorem, we will track the ratio of weights of the attention scores $\sos$'s with  to the highest attention score $\sosi{0}$ . Let us denote the ratio of the attention scores as follows: 
\begin{align*}
    \qqi{i}(t) = \frac{\sosi{i}(t)}{\sosi{0}(t)} = \exp{\left( \aai{i}(t) - \aai{0}(t) \right)}.    
\end{align*}
Using the $\qqi{i}$'s the softmax weights can be written as,
\begin{align*}
\sosi{i} &= \frac{\qqi{i}}{1 + \sum_{j = 1}^{d-1} \qqi{j}} = \frac{\qqi{i}}{\cQ},  \\
\end{align*}  
where
\begin{align*}
    \cQ &= 1 + \sum_{j = 1}^{d-1} \qqi{j}. 
\end{align*}


\paragraph{Evolution of the ratios.} To compute the evolution of $\qqi{i}$, we will first compute the following dynamics,

\begin{align*}
    d {e^{-\aai{i}}} &= - \frac{\gamma}{\cA} \left( \uui{i} - \scal{\uu}{\sos} \right) d{t}, \\
    d { \left( \uui{i} - \uui{j} \right) } &= \frac{\gamma}{\cA} \left( e^{\aai{i}} - e^{\aai{j}} \right) d{t}, \\
    d { \left( e^{-\aai{i}} - e^{-\aai{j}} \right) } &= - \frac{\gamma}{\cA} \left( \uui{i} - \uui{j} \right) d{t}. \\
    d { \, e^{\aai{i}}} &= \gamma \frac{\left(e^{\aai{i}}\right)^2}{\cA} \left( \uui{i} - \scal{\uu}{\sos} \right) d{t}.
\end{align*}
Lets compute, 
\begin{align*}
    d \, \qqi{i} =  d { \left( e^{ \aai{i}}  e^{-\aai{0}}\right) } &= e^{-\aai{0}} d {e^{\aai{i}}} + e^{\aai{i}} d {e^{-\aai{0}}}, \\     
    &= \gamma \frac{e^{-\aai{0}} \left(e^{\aai{i}}\right)^2}{\cA} \left( \uui{i} - \scal{\uu}{\sos} \right) d{t} - \gamma \frac{e^{\aai{i}}}{\cA} \left( \uui{0} - \scal{\uu}{\sos} \right) d{t}, \\
    &= \gamma \frac{e^{\aai{i}}}{\cA} \left[ e^{-\aai{0}} e^{\aai{i}} \left( \uui{i} - \scal{\uu}{\sos} \right) - \left( \uui{0} - \scal{\uu}{\sos} \right) \right] d{t}, \\
    &= \gamma \sosi{i} \left[ \,  \qqi{i} \left( \uui{i} - \scal{\uu}{\sos} \right) - \left( \uui{0} - \scal{\uu}{\sos} \right) \right] d{t}, \\
    &= \gamma \, \, \frac{\qqi{i}}{\cQ} \, \, \left[ \,  \qqi{i} \left( \uui{i} - \scal{\uu}{\sos} \right) - \left( \uui{0} - \scal{\uu}{\sos} \right) \right] d{t}.
\end{align*}

% \begin{align*}
% \qqi{i} &=  e^{\aai{i} - \aai{0}}, \\
% \sosi{i} &= \frac{\qqi{i}}{1 + \sum_{j = 1}^{d-1} \qqi{j}} = \frac{\qqi{i}}{\cQ},  \\
% \text{where } \cQ &= 1 + \sum_{j = 1}^{d-1} \qqi{j}. 
% \end{align*} 
Using the above derivation, the dynamics of $\qqi{i}$ can be written as,
\begin{align*}
    \der{\qqi{i}} &= \gamma \frac{\qqi{i}}{\cQ} \left[ \qqi{i} \left( \uui{i} - \scal{\uu}{\sos} \right) - \left( \uui{0} - \scal{\uu}{\sos} \right) \right] \der{t}, \\
    &= \gamma \frac{\qqi{i}}{\cQ} \left[ - \qqi{i} \left( \uui{0} - \uui{i} \right) - \left(  1 - \qqi{i} \right)\left(  \uui{0} - \scal{\uu}{\sos} \right) \right] \der{t}.
\end{align*}

From the previous therorem all time $t \geq 0$, the following properties hold along the trajectory of the gradient flow, 
\begin{enumerate}[label=(\alph*)]
    \item For all $i = 1, 2, \ldots, p-1$, $\, \, 0 \leq \qqi{i}(t) \leq 1$  and therefore $ 1 \leq \cQ \leq p$ for all time $t \geq 0$.
    \item For all $i = 1, 2, \ldots, p-1$, $\uui{0}(t) - \uui{i}(t) \geq \delta $ for all time $t \geq 0$.
\end{enumerate}

Using these properties, we will have some asymptotic bounds on relevant quantities. First, the evolution of $\qqi{i}$ can be bounded as follows,
\begin{align*}
    \der{\qqi{i}} &\leq -  \frac{\gamma}{\cQ} \left[ \qqi{i}^2 + \qqi{i} (1 - \qqi{i}) ( 1 - \frac{1}{\cQ}) \right] \delta \der{t}, \\
    \der{\qqi{i}} &\leq - \frac{\gamma}{\cQ} \left[ \qqi{i} - \frac{\qqi{i}(1-\qqi{i})}{\cQ} \right] \delta \der{t} , \\
    &\leq -\frac{\gamma}{\cQ} \left[ \qqi{i} - \qqi{i}(1 - \qqi{i}) \right] \ delta \der{t}, \\
    &\leq - \frac{\gamma}{\cQ} \qqi{i}^2 \delta \der{t} \leq -  \frac{\gamma}{p} \qqi{i}^2 \delta \der{t}.
\end{align*} 
Using this, we can have the following bound on $\qqi{i}$,
\begin{align*}
    \int_{1}^{\qqi{i}(t)} \frac{1}{q^2} \der{q} &\leq - \frac{\delta}{p} \int_{0}^{t} \gamma(h) \der{h}, \\
    - \left( \frac{1}{\qqi{i}(t)} - 1 \right) &\leq - \frac{\delta}{p} \int_{0}^{t} \gamma(h) \der{h}, \\
    \qqi{i}(t) &\leq \frac{1}{1 + \frac{\delta}{p} \int_{0}^{t} \gamma(h) \der{h}}.
\end{align*}
Note that this bound is irrespective of $i$. The goal of the rest of the proof is to show that the integral $\int_{0}^{t} \gamma(h) \der{h}$ diverges as $t \to \infty$.
Recall that,
\begin{align*}
    \gamma = \frac{1}{1 + \exp\left(  \scal{\bbeta_*}{\bbeta}  \right)} = \frac{1}{1 + \exp\left(  \sum_{j} \sosi{j} \uui{j}  \right)}. 
\end{align*} 
In the above expression, $\uui{j} \leq \uui{0}$ for all $j$.
\begin{align*}
    \uui{j} \leq \uui{0} &\implies \sum_{j} \sosi{j} \uui{j} \leq \uui{0} \sum_{j} \sosi{j} = \uui{0}, \\
    &\implies 1 + \exp\left(  \sum_{j} \sosi{j} \uui{j}  \right) \leq 1 + \exp\left(  \uui{0}  \right), \\
    &\implies \gamma \geq \frac{1}{1 + \exp\left(  \uui{0}  \right)}.
\end{align*}
First, lets check evolution of $\uui{0}$, 
\begin{align*}
    d{\uui{0}} &=  \gamma  \sosi{0} dt =  \gamma \frac{1}{\cQ} d{t} \geq \frac{\gamma}{p} d{t} \geq \frac{1}{p(1 + \exp( \uui{0} ))} d{t}.
\end{align*}
From this evolution we have the following, 
\begin{align*}
    \int_{0}^{t} d h &\leq p \int_{\uui{0}(0)}^{\uui{0}(t)} (1 + \exp(z)) \der{z}, \\
    t &\leq p \left[ \uui{0}(t) - \uui{0}(0) + \exp(\uui{0}(t)) - \exp(\uui{0}(0)) \right], \\
    &\leq p \left[ \uui{0}(t) + \exp(\uui{0}(t)) + c_0 \right], \\
    &\leq p \left[ 2 \exp(\uui{0}(t)) + c_0 \right], \\
    &\implies \exp(\uui{0}(t)) \geq \frac{t}{2p} - c_0, \\
    &\implies \uui{0}(t) \geq \log\left(  \frac{t}{2p} - c_0  \right).
\end{align*}
We have the following bound on $\uui{0}(t)$,
\begin{align} \label{eq:u0-bound}
    \uui{0}(t) &\geq \log\left(  \frac{t}{2p} - c_0  \right).
\end{align}
We have established a lower bound on $\uui{0}(t)$. We will use this to lower bound the integral $\int_{0}^{t} \gamma(h) \der{h}$. We have, 
\begin{align*}
\der{\uui{0}} &= \frac{\gamma}{\cQ} \der{t} \leq \gamma \der{t}, 
\end{align*}
This implies,
\begin{align*}
 \uui{0} - \uui{0}(0) \leq \int_{0}^{t} \gamma(h) \der{h} &\implies \int_{0}^{t} \gamma(h) \der{h} \geq \uui{0}(t) - \uui{0}(0).
\end{align*}
Combining this we have the following, 
\begin{align} \label{eq:gamma-t-bound}
\int_{0}^{t} \gamma(h) \der{h} &\implies \int_{0}^{t} \gamma(h) \der{h} \geq \uui{0}(t) - \uui{0}(0) \geq \log\left(  \frac{t}{2p} - c_0  \right) - \uui{0}(0).
\end{align}
This implies that 
\begin{align}\label{eq:gamma-diverges}
    \lim_{t \to \infty} \int_{0}^{t} \gamma(h) \der{h} &\to \infty.
\end{align}
Using the bound on $\qqi{i}(t)$ derived earlier, we have,
\begin{align*}
    \qqi{i}(t) &\leq \frac{1}{1 + \frac{\delta}{p} \int_{0}^{t} \gamma(h) \der{h}} \leq \frac{p}{p + \delta \left( \log\left( t\right) - c_0  \right)}.
\end{align*}
Using the divergence property from equation~\eqref{eq:gamma-diverges}, we have,
\begin{align*}
    \lim_{t \to \infty} \qqi{i}(t) &\to 0.
\end{align*}

The only part that is left to be proved is convergence of the loss to zero. 
Note that 
\begin{align*}
    \scal{\bbeta_*}{\bbeta(t)} &= \sum_{j} \sosi{j}(t) \uui{j}(t) \leq \uui{0}(t) \sum_{j} \sosi{j}(t) \geq \frac{1}{p} \uui{0}(t).
\end{align*}
and we have already established that $\uui{0}(t) \to \infty$ as $t \to \infty$. Therefore, we have,
\begin{align*}
    \lim_{t \to \infty} \scal{\bbeta_*}{\bbeta(t)} &\to \infty.
\end{align*}
Using this in the expression of the logistic loss function in Eq.~\eqref{eq:logistic_beta}, we have,
\begin{align*}
    \lim_{t \to \infty} \bloss(\bbeta(t)) &= \lim_{t \to \infty} \log \left( 1 + \exp\left( -  \scal{\bbeta_*}{\bbeta(t)}  \right) \right) = 0.
\end{align*}    
\end{proof}

\begin{lemma}~\label{lem:log-non-maximal-rates} We will make two additional remarks regarding the theorem above which gives asymptotic bounds on the non-maximal scores and projection parameters. 
\begin{enumerate}[label=(\alph*)]
    \item The projection co-ordinates $\uui{i}(t)$ for $i \neq 0$ does not grow unbounded with time, i.e., there exists a constant $C > 0$ such that for all time $t \geq 0$, we have $\uui{i}(t) \leq C$. As $\uui{0} \to \infty$, this implies that $\VV$ is close to a rank 1 structure. 
    \item The attention scores for $i \neq 0$ satisy $\sosi{i} = O(\frac{1}{\log^2 (t)})$. 
\end{enumerate}
\end{lemma}    

\begin{proof}
From Eq.~\eqref{eq:gamma-t-bound}, we can say that $\gamma = O(1/t)$ and from the evolution of $\uui{0} - \uui{1}$, we have that, 
\begin{align*}
    d \{\uui{0} - \uui{i}\} &= \frac{\gamma}{\cQ} ( 1 - \qqi{i} ) \der{t} 
\end{align*}
As $\qqi{i} \to 0$, there exists a time $t_0$ after which we have $1 - \qqi{i} \geq 1/2$. Therefore, for $t \geq t_0$, we have,
\begin{align*}
    d \{\uui{0} - \uui{i}\} &\geq \frac{\gamma}{2p} \der{t} \\
    \uui{0}(t) - \uui{i}(t) &\geq \frac{1}{2p} \int_{t_0}^{t} \gamma(h) \der{h} + c_1 = \Omega(\log(t)),
\end{align*}

Using this in the evolution of $\qqi{i}$, we have,
\begin{align*}
    \der{\qqi{i}} &\leq - \frac{\gamma}{\cQ} \qqi{i}^2 \delta \der{t} \leq -  \frac{\gamma}{p} \qqi{i}^2 \delta \der{t}.
\end{align*}

    From the above equations~\eqref{eq:u0-bound} we have that $\uui{0}(t) \geq \log\left(  \frac{t}{2p} - c_0  \right).$ and~\eqref{eq:gamma-t-bound}, we have the following bound on $\qqi{i}(t)$,
\begin{align*}
    \der{\qqi{i}}
    &= \gamma \frac{\qqi{i}}{\cQ} \left[ - \qqi{i} \left( \uui{0} - \uui{i} \right) - \left(  1 - \qqi{i} \right)\left(  \uui{0} - \scal{\uu}{\sos} \right) \right] \der{t}. \\
    &\geq - c \, \frac{\gamma}{p} \, \qqi{i}^2 \, \log{t} \geq - c  \, \qqi{i}^2 \, \frac{\log{t}}{t}.
\end{align*}
using the asymptotic bound on $\gamma, \uui{0} - \uui{i}$. Integrating this we get that 
\begin{align*}
    \int_{1}^{\qqi{i}(t)} \frac{1}{q^2} \der{q} &\geq - c \int_{1}^{t} \frac{\log{h}}{h} \der{h}, \\
    - \left( \frac{1}{\qqi{i}(t)} - 1 \right) &\geq - c \left( \frac{(\log{t})^2}{2} - \frac{(\log{1})^2}{2} \right), \\
    \qqi{i}(t) &\leq \frac{1}{1 + c \, \frac{(\log{t})^2}{2}} = O\left( \frac{1}{\log^2(t)} \right).
\end{align*}

Now coming to the evolution of $\uui{i}$, we have,
\begin{align*}
    d{\uui{i}} &=  \gamma  \sosi{i} dt =  \gamma \frac{\qqi{i}}{\cQ} d{t} \leq \gamma \, \qqi{i} \, d{t} \leq \, c \frac{1}{t \log^2(t)} d{t}.   
\end{align*}
Integrating again, we get, 
\begin{align*}
    \uui{i}(t) - \uui{i}(t_0) &\leq c \int_{t_0}^{t} \frac{1}{h \log^2(h)} d{h}, \\
    &\leq c \left[ - \frac{1}{\log(h)} \right]_{t_0}^{t}, \\
    &\leq c \left( \frac{1}{\log(t_0)} - \frac{1}{\log(t)} \right) \leq c'.
\end{align*}
This proves the remark.
\end{proof}

\section{Linear Regression}

Here we consider the case of loss function given by the squared loss, i.e., 
\begin{align} \label{eq:sq_loss_beta}
    \bloss(\bbeta) = \| \bbeta - \bbeta_* \|^2.
\end{align}


The gradient flow dynamics for the parameters $\vv$ and $\aa$ can be written as,
\begin{align}\label{eq:gradient-flow-v-a-square}
    \mdot{\VV} &= ( \bbeta_* - \bbeta ) \sos^{\top}, \\
    \mdot{\aa} &= \left[ \diag{\sos} - \sos \sos^{\top} \right] \VV^{\top} ( \bbeta_* - \bbeta ).
\end{align}

Let $\beta_{\perp}$ be any direction be orthogonal to $\bbeta_*$. We can decompose $\vv$ as,
\begin{align*}
   \der{ \left( \beta_{\perp} \right)^{\top} \VV }  &= \beta_{\perp}^{\top} \left( \bbeta_* - \VV \sos \right) \sos^{\top} \der{t} = -  \left( \beta_{\perp} \right)^{\top} \VV \sos \sos^{\top} \der{t}. \\
\end{align*}


\Regression*
\begin{proof}
    Recall that the gradient flow dynamics of the parameters $\VV$ and $\aa$ are given by,
    \begin{align}
    \mdot{\VV} &= ( \bbeta_* - \bbeta ) \sos^{\top}, \\
    \mdot{\aa} &= \left[ \diag{\sos} - \sos \sos^{\top} \right] \VV^{\top} ( \bbeta_* - \bbeta ).
\end{align}

Let $\beta_{\perp}$ be any direction be orthogonal to $\bbeta_*$. We can decompose $\vv$ as,
\begin{align*}
   \der{ \left( \beta_{\perp} \right)^{\top} \VV }  &= \beta_{\perp}^{\top} \left( \bbeta_* - \VV \sos \right) \sos^{\top} \der{t} = -  \left( \beta_{\perp} \right)^{\top} \VV \sos \sos^{\top} \der{t}.
\end{align*}

Under this assumption, we know that the initial projection of $\VV(0)$ along $\beta_{\perp}$ is zero, i.e., $\left( \beta_{\perp} \right)^{\top} \VV(0) = 0$. Therefore, from the above equation, we have,
\begin{align*}
    \left( \beta_{\perp} \right)^{\top} \VV &= 0, \quad \forall t \geq 0.
\end{align*}
This implies that $\VV$ is rank one and $\bbeta_*$ is a left singular vector of $\VV$ and $$ \VV = \frac{\bbeta_*}{\nor{\bbeta_*}^2} \uu^{\top}. $$
Using this is the gradient flow dynamics of $\aa$, we have,
\begin{align*}
    \mdot{\aa} &= \left[ \diag{\sos} - \sos \sos^{\top} \right] \frac{\uu}{\nor{\bbeta_*}^2} \left( \bbeta_*^{\top} \bbeta_* - \bbeta_*^{\top} \bbeta \right), \\
    &= \left[ \diag{\sos} - \sos \sos^{\top} \right] \uu \left( 1 - \scal{\bbeta_*}{\bbeta} / \nor{\bbeta_*}^2 \right).     
\end{align*}
Similarly, the gradient dynamics of $\uu$ can be written as,
\begin{align*}
    \mdot{\uu} &= \nor{\bbeta_*}^2 \left( 1 - \scal{\bbeta_*}{\bbeta} / \nor{\bbeta_*}^2 \right) \sos.
\end{align*}
This finishes the proof. 
\end{proof}


\subsection{Remarks on the linear regression case}
Note that the dynamics of $\aa$ and $\uu$ are similar to the dynamics in the classification case with logistic loss. The only difference is the scaling factor $\left( 1 - \scal{\bbeta_*}{\bbeta} / \nor{\bbeta_*}^2 \right)$ which depends on the current loss. Theorem~\ref{thm:no-crossing-logitstic} implies the no-crossing property of the projections and attention scores in this case as well in the case of linear regression.

Now we will turn to the strength of the polarization, i.e., $\int \gamma(s) ds .$ Following the same approach as in Theorem~\ref{thm:sparsity-logistic}, we can establish the following result regarding the sparsity of the attention scores in the case of linear regression. From now on, assume that $\nor{\bbeta_*} = 1$ for brevity. From the Lemma~\ref{lem:convergence-general-loss}, we have that,
\begin{align*}
    \frac{ d \ell(\bbeta)}{dt} &\leq - \frac{1}{p} \nor{\nabla_{\bbeta} \ell}^2.
 \end{align*}
Using this in the context of linear regression, we get, 
\begin{align*}
    \ell(\bbeta) &= \frac{1}{2} \| \bbeta - \bbeta_* \|^2, \\
    \nabla_{\bbeta} \ell &= ( \bbeta - \bbeta_* ), 
\end{align*}
Substituting these in the convergence lemma, we have, 
\begin{align*}
     \frac{d \| \bbeta - \bbeta_* \|^2 }{dt} &\leq - \frac{1}{p} \| \bbeta - \bbeta_* \|^2. 
\end{align*}
Hence the loss converges to zero exponentially fast and note that 
$\gamma = \scal{ \bbeta_*}{\bbeta_* - \bbeta} \leq \nor{\bbeta_*} \nor{\bbeta_* - \bbeta}$. Hence that $\gamma$ also converges to zero exponentially fast. Therefore, the integral $\int_{0}^{t} \gamma(h) \der{h}$ converges to a finite value as $t \to \infty$. This implies that the attention scores do not converge to a one-hot vector in the case of linear regression.



\section{Supporting Material}

\subsection{Some properties of the value-softmax model.}

\begin{lemma}
\textbf{Invariance.} The loss function have the following invariance property:    

\begin{enumerate}[label=(\alph*)]
    \item For any $c \in \R$, we have, 
    \begin{align*}
        \loss(\VV, \aa) = \loss(\VV, \aa + c \one).
    \end{align*}
    and taking the gradient with respect to $c$ we get the following invariance in the gradients: 
    \begin{align*}
        \scal{\nabla_{\aa} \loss(\VV, \aa)}{\one} = 0.      
    \end{align*}
    \item The gradient flow dynamics has the following invariance:  
    \begin{align*}
        \sum_{i=1}^{d} \frac{d \aai{i}}{dt} = 0 \implies \sum_{i=1}^{d} \aai{i}(t) = \sum_{i=1}^{d} \aai{i}(0), \quad \forall t \geq 0.
    \end{align*}
\end{enumerate}
\end{lemma}


\begin{lemma}\label{lem:convergence-general-loss}\textbf{Convergence.}
Consider the gradient flow dynamics given by the equations~\eqref{eq:gradients-general-loss} over the loss function $\loss(\VV, \aa)$, 
\begin{align*}
    \frac{d \, \bloss(\bbeta)}{dt} &\leq - \frac{1}{d} \nor{\nabla_{\bbeta} \bloss}^2 
\end{align*} 
\end{lemma}
\begin{proof}
    With $\rrr = - \nabla_{\bbeta} \bloss $, the time derivatives of the parameters are given by, 
    \begin{align*}
        \frac{d\VV}{dt} &=  \rrr \sos^{\top},  \\
        \frac{d\sos}{dt} &=  \left( \diag{\sos} - \sos \sos^{\top} \right)^2 \VV^{\top} \rrr.
    \end{align*}
    Considering the evolution of $\bbeta = \VV \sos$, we have,
    \begin{align*}
        \frac{d \bbeta}{dt} &= \frac{d \VV \sos}{dt} = \frac{d\VV}{dt} \sos + \VV \frac{d\sos}{dt}, \\
        &= \left( \rrr \sos^{\top} \right) \sos + \VV \left( \diag{\sos} - \sos \sos^{\top} \right)^2 \VV^{\top} \rrr, \\
        &= \left[ \nor{\sos}^2 \I +  \VV \left( \diag{\sos} - \sos \sos^{\top} \right) \VV^{\top} \right] \rrr. 
    \end{align*}
    Let us denote the matrix in the brackets as $\MM(\sos, \VV)$, i.e., 
    \begin{align*}
        \MM(\sos, \VV) = \nor{\sos}^2 \I +  \VV \left( \diag{\sos} - \sos \sos^{\top} \right) \VV^{\top}.
    \end{align*}    
    Using this notation, we can write the evolution of $\bbeta$ as follows:
    \begin{align*}
        \frac{d \bbeta}{dt} &=  \MM(\sos, \VV) \rrr = - \MM(\sos, \VV) \nabla_{\bbeta} \bloss. 
    \end{align*}
    Considering the time derivative of the loss function $\bloss(\bbeta)$, we have,
    \begin{align*}
        \frac{d \bloss(\bbeta)}{dt} &= \scal{\nabla_{\bbeta} \bloss}{\frac{d \bbeta}{dt}} = - \scal{\nabla_{\bbeta} \bloss \, \, }{\, \, \MM(\sos, \VV) \, \nabla_{\bbeta} \bloss}. 
    \end{align*}
    Now we will show that the matrix $\MM(\sos, \VV)$ is positive definite. First, 
    \begin{align*}
        \nor{\sos}^2  \geq \sum_{i=1}^{d} \sosi{i}^2 \geq \frac{1}{d} \left( \sum_{i=1}^{d} \sosi{i} \right)^2 = \frac{1}{d}.
    \end{align*}
    The next term is positive semi-definite since for any vector $\uu \in \R^{d}$, we have,
    \begin{align*}
        \uu^{\top} \VV \left( \diag{\sos} - \sos \sos^{\top} \right) \VV^{\top} \uu &= \left( \VV^{\top} \uu \right)^{\top} \left( \diag{\sos} - \sos \sos^{\top} \right) \left( \VV^{\top} \uu \right), \\
        &= \mu^{\top} \left( \diag{\sos} - \sos \sos^{\top} \right) \mu, \quad \text{ where } \mu = \VV^{\top} \uu, \\
        &= \sum_{i=1}^{d} \sosi{i} \mu_i^2 - \left( \sum_{i=1}^{d} \sosi{i} \mu_i \right)^2 \geq 0,
    \end{align*}
    where the last inequality follows from the Cauchy-Schwarz inequality. Hence we have,
    \begin{align*}
        \uu^{\top} \MM(\sos, \VV) \uu \geq \frac{1}{d} \nor{\uu}^2, \quad \forall \, \uu \in \R^{d}.
    \end{align*}
    This shows that the matrix $\MM(\sos, \VV)$ is positive definite. Using the positive definiteness of $\MM(\sos, \VV)$, we have,
    \begin{align*}
        \frac{d \bloss(\bbeta)}{dt} &\leq - \frac{1}{d} \nor{\nabla_{\bbeta} \bloss}^2 \leq 0.
    \end{align*}
    This shows that the loss function $\bloss(\bbeta)$ is non-increasing along the gradient flow dynamics of $\VV$ and $\sos$.
\end{proof}


\begin{property} \label{prop:jacobians}
For the value-softmax model, i.e., $\bbeta = \VV \sf(\aa)$, this property summarizes the jacobians of the different components.
    \begin{itemize}
        \item The jacobian of the softmax function $\sf: \R^{d} \to \Delta^{d}$ :
        \begin{align}\label{eq:softmax_jacobian}
            \jacb{\aa}{\sf(\aa)} = \diag{\sf(\aa)} - \sf(\aa) \sf(\aa)^{\top}. 
        \end{align}
        With $\sos = \sf(\aa)$, the jacobian can be written as
        \begin{align}\label{eq:s_jacobian}
            \jacb{\aa}{\sos} = \diag{\sos} - \sos \sos^{\top}. 
        \end{align}
        \item The jacobian with respect to the $\bbeta$ with respect to $\VV$ and $\aa$ are given by
        \begin{align}
            \jacb{\VV}{\bbeta} &= \sos^{\top} \otimes \I_{d}, \\
            \jacb{\sos}{\bbeta} &= \VV, \\
            \jacb{\aa}{\bbeta} &= \VV \jacb{\aa}{\sos} = \VV \left( \diag{\sos} - \sos \sos^{\top} \right).
        \end{align}
        It is needed to be clarified that the jacobian $\jacb{\VV}{\bbeta}$ is a matrix of size $d \times d^2$ and it is constructed in the following way:
        \begin{align*}
            \jacb{\VV}{\bbeta} = \begin{bmatrix}
                \jacb{\vvc{1}}{\bbeta} & \jacb{\vvc{2}}{\bbeta} & \ldots & \jacb{\vvc{d}}{\bbeta}
            \end{bmatrix} = \sos^{\top} \otimes \I_{d}.
        \end{align*}
        where each row is $ \jacb{\mathrm{vec}{(\VV)}}{\bbeta_i} $ . 
    \end{itemize}
\end{property}

\subsection{Normalization} 
Here, we will consider the case of general normalization functions and establish that the no-crossing property holds for the logistic loss as well as the regression loss function . In particular, 
consider a general monotonically increasing differentiable function $f$, and define the normalization $\sf_{f}: \R^{p} \to \R^{p}$ for any $a$ in $\R^{p}$ as
\begin{align*}
  \sf_{f}(a)_i = \frac{f(a_i)}{\sum\limits_{j=1}^{p} f(a_j)}.
\end{align*}
For the logistic loss function, the gradient flow dynamics is given by
\begin{align} \label{eq:general-normalization-logistic} 
  \der{\uu} &=  \gamma(\bbeta) \, \|\bbeta_*\|^2   \sos_{f} \, \der{t}, \\
  \der{\aa} &=  \gamma(\bbeta) \, \frac{f'(\aa)}{\sum\limits_{j=1}^{p} f(\aa_j)} \odot ( \uu - \scal{\uu}{\sf_f} \one ) \, \der{t}.
\end{align}


\begin{lemma}
    Consider the gradient flow dynamics given by the equation Eq.~\eqref{eq:general-normalization-logistic} over the logistic loss function $\bloss(\bbeta)$. For the initialization given by the Assumption~\ref{ass:init}, the following no-crossing property holds for all time $t \geq 0$,
    \begin{align*}
        \uui{0}(t) > \uui{1}(t) > \uui{2}(t) \ldots > \uui{p-1}(t), \\
        \aai{0} (t) \geq  \aai{1} (t) \geq \aai{2} (t) \ldots \geq \aai{p-1} (t).
     \end{align*}    
\end{lemma}
\begin{proof}
    Let $\bF = \sum_{j} f(\aa_j)$ and let $$ G(s) = \int_{0}^{s} \frac{1}{f'(v)} dv. $$ Using this to we rewrite the co-ordinate wise dynamics as 
    \begin{align*}
        d\uui{i} &= \frac{\gamma}{\bF} f(\aai{i}) dt, \\ 
        d\aai{i} &= \frac{\gamma}{\bF} f'(\aai{i}) \left( \uui{i} - \scal{\uu}{\sos_f} \right) dt, \\
        \frac{d \aai{i}}{f'(\aai{i})} &= \frac{\gamma}{\bF} \left( \uui{i} - \scal{\uu}{\sf_f} \right) dt. \\
        d{ G(a_i) } &= \frac{\gamma}{\bF} \left( \uui{i} - \scal{\uu}{\sf_f} \right) dt, \\
    \end{align*}
For any $i,j$, taking the differences, 
    \begin{align*}
        d \left( G(\aai{i}) - G(\aai{j}) \right) &= \frac{\gamma}{\bF} \left( \uui{i} - \uui{j} \right) dt. \\
        d \left( \uui{i} - \uui{j} \right) &= \frac{\gamma}{\bF} \left( f(\aai{i}) - f(\aai{j}) \right) dt.
    \end{align*}

    We can take a potential function $\Phi_{ij} = \left( G(\aai{i}) - G(\aai{j}) \right)\left( \uui{i} - \uui{j} \right)$, using this we have, 
    \begin{align*}
        d \Phi_{ij} &= \left( \uui{i} - \uui{j} \right) d \left( G(\aai{i}) - G(\aai{j}) \right) + \left( G(\aai{i}) - G(\aai{j}) \right) d \left( \uui{i} - \uui{j} \right), \\
        &= \frac{\gamma}{\bF} \left[ \left( \uui{i} - \uui{j} \right)^2 + \left( G(\aai{i}) - G(\aai{j}) \right) \left( f(\aai{i}) - f(\aai{j}) \right) \right] dt. 
    \end{align*}
    Note that $f$ increases monotonically, hence $f'(a) > 0$ everywhere, which implies $1/f'(a) > 0$ everywhere and G also increases monotonically. We assume that it does not hit $0$ in finite time, so the term $\left( G(\aai{i}) - G(\aai{j}) \right) \left( f(\aai{i}) - f(\aai{j}) \right)$ is always non-negative. This implies that $d \Phi_{ij} \geq 0$, ensuring that the potential function $\Phi_{ij}$ does not decrease over time. Consequently, the no-crossing property holds for all time $t \geq 0$.
\end{proof}

% \subsection{Experiments with Value-Softmax Model.}


% \section{}


% \section{}



\section{Additional Experimental Details}
\label{app:experimental-details}

\subsection{Induction heads acting as attention sinks}

\paragraph{Task format.} The structure of each object in the dataset is illustrated in Figure \ref{fig:induction_task_format}. Each object consists of a beginning-of-sequence token, $k$ randomly sampled tokens from the vocabulary $t_1,\dots,t_k$, a token $t_i'$ uniquely corresponding to a randomly chosen $t_i$, followed by $t_{i+1}$. We experimented both with $t_i' = t_i$ and $t_i'$ being a different symbol, observing that the latter typically leads to less attention sinks forming.

The loss is calculated as $L_\mathrm{total} = \lambda \cdot (1/k) L_\mathrm{random} + L_\mathrm{induction}$. We tested $\lambda = 0$ and $\lambda = 0.2$. In the former case, the model is only penalized on incorrect predictions of the last token and not incentivized to keep a constant uniform prediction on random tokens. This leads to all heads duplicating each other and no almost no attention sinks for any attention type.

\begin{figure}[h]
  \centering
  \includegraphics[width=0.5\textwidth]{figs/induction.pdf}
  %\vspace{-7pt}
  \caption{Task format in the induction heads experiment. The arrows illustrate the induction head mechanism.}
  \label{fig:induction_task_format}
\end{figure}

\begin{table}[h]
  \centering
  \renewcommand{\arraystretch}{1.15}
  \setlength{\tabcolsep}{10pt}
  \begin{tabular}{l c l c}
    \hline
    \textbf{Hyperparameter} & \textbf{Value} & \textbf{Hyperparameter} & \textbf{Value} \\
    \hline
    Layers           & 2   & Epochs     & 10 \\
    Hidden dim       & 128 & Train size         & 100{,}000 \\
    Feedforward dim  & 512 & Eval size          & 1{,}000 \\
    Vocab size       & $51^*$  & Sequence length & 13 \\
    Causal           & \texttt{true} & Batch size     & 256 \\
    Use LayerNorm    & \texttt{true} & Use MLP        & \texttt{true} \\
    \hline
    Heads            & $\bf \{1,2,4,8\}$ & $t_i' = t_i$ & $\{\texttt{true},\textbf{\texttt{false}}\}$ \\
    Learning rate    & $\bf \{10^{-2},10^{-3},10^{-4}\}$ & $\lambda$ & $\{0.0,\bf{0.2}\}$ \\
    Attention type   & all in Table \ref{table:attention-types} & Positional type  & \{\texttt{concat}, \textbf{\texttt{additive}}, \texttt{no}\} \\
    Seed             & $\bf \{0,1,2,3,4\}$ & & \\
    \hline
  \end{tabular}
  \vspace{3pt}
  \caption{Configuration and hyperparameters for the induction experiment. The bold values indicate the setup shown in Figures \ref{fig1} and \ref{fig:main-exp-induction}. Results for other setups are additionally presented in Appendix \ref{app:results-induction} Vocabulary size is increased by 50 if $t_i' \ne t_i$.}
  \label{tab:induction-config-sweep}
\end{table}

\paragraph{Hyperparameters.}

In this experiment, the model is a standard 2-layer $H$-head causal Transformer with pre-normalization, similar to GPT-2 \cite{radford2019language}. The full hyperparameter and configuration details are available in Table \ref{tab:induction-config-sweep}.

\paragraph{Metrics.}

We discard all runs that did not achieve induction accuracy above $95\%$. A sink score for each head is measured as an average proportion of attention weight assigned to the BOS token on the positions where no informed prediction is possible. We consider a head \emph{a sink head} if it's sink score is higher than 0.9.
The computation can be expressed via the following pseudocode:
{\small \begin{verbatim}
    A = attention_matrix  # [layer, head, sample, query_seq_len, key_seq_len]

    # index 0 is removed because it's BOS. 
    # indices -1 and -2 are removed because prediction is not random.
    bos_score[L, H] = A[L, H, :, 1 : -2, 0]
    total_weight[L, H] = A[L, H, :, 1 : -2, :].sum(axis=-1)

    # clipping affects attention types without positivity. doesn't change is_sink.
    sink_score[L, H] = (bos_score[L, H] / total_weight[L, H]).mean().clip(0, 1) 
    is_sink[L, H] = (sink_score[L, H] > 0.9)
\end{verbatim}}

We only plot the results from the second layer, as all heads in the first layer typically attend to the previous token.

\subsection{Difference in attention sparsity between pretrained softmax and sigmoid LLMs}

We use the weights and the codebase\footnote{https://github.com/apple/ml-sigmoid-attention} for the pretrained 7B softmax and sigmoid models released by \citet{ramapuramtheory}. We run both models on 1000 samples from a subset\footnote{https://huggingface.co/datasets/NeelNanda/pile-10k} of the Pile dataset \citep{gao2020pile} with the maximum sequence length of 256. We measure the sparsity score of each head as the average proportion of total attention weight assigned to the max-logit token. In pseudocode,
{\small \begin{verbatim}
    A = attention_matrix  # [layer, head, sample, query_seq_len, key_seq_len]
    max_attn_weight = A.max(axis=-1)
    total_attn_weight = A.sum(axis=-1)
    sparsity_score[L, H] = (max_attn_weight / total_attn_weight).mean()
\end{verbatim}}

In Figure \ref{fig:main-llm-pretrained}, we plot the distribution of sparsity scores for all heads in softmax and sigmoid models.

\subsection{Emerging imbalance in token influence on classifier's predictions}

\paragraph{Task format.} The task is sequence classification into $K$ labels. The vocabulary is split into $K$ disjoint groups $V^{(1)},\dots,V^{(K)}$. Each object with a label $k$ is a random sequence of tokens from the $k$-th group: $x_i = v^{(k_i)}_{j_{i1}},\dots,v^{(k_i)}_{j_{in}}, y_i=k$. \texttt{CLS} token is appended to the input. Knowledge of any of the input tokens is sufficient to achieve perfect in-distribution accuracy. The task is illustrated in Figure \ref{fig:cls_task_format} (left).

\begin{figure}[h]
  \centering
  \includegraphics[width=1.0\textwidth]{figs/classification.pdf}
  \caption{Task format in the classification experiment and the token swap procedure. Cell brightness illustrates the attention weight assigned to each token.}
  \label{fig:cls_task_format}
\end{figure}

\begin{table}[h]
  \centering
  \renewcommand{\arraystretch}{1.15}
  \setlength{\tabcolsep}{10pt}
  \begin{tabular}{l c l c}
    \hline
    \textbf{Hyperparameter} & \textbf{Value} & \textbf{Hyperparameter} & \textbf{Value} \\
    \hline
    Layers           & $\{{\bf1},2\}$ & Epochs         & 10 \\
    Hidden dim       & 128          & Train size     & 10{,}000 \\
    Feedforward dim  & 512          & Eval size      & 1{,}000 \\
    Vocab size       & $\{\bf{41},81,161\}$ & Sequence length & 11 \\
    Labels & 4     & Causal        & \texttt{false} \\
    Batch size       & 256          & Positional type & \texttt{no} \\
    \hline
    Heads            & $\{{\bf 1},2,4\}$ & Use MLP        & $\{\textbf{\texttt{true}},\texttt{false}\}$ \\
    Use LayerNorm    & $\{\textbf{\texttt{true}},\texttt{false}\}$ & Learning rate & $\bf\{10^{-2},10^{-3},10^{-4}\}$ \\
    Attention type   & all in Table \ref{table:attention-types} & Seed & $\bf\{0,1,2,3,4\}$ \\
    \hline
  \end{tabular}
  \vspace{3pt}
  \caption{Configuration and hyperparameters for the classification experiment. The bold values indicate parameters used for the setup in Figures \ref{fig:main-exp-classification} and \ref{fig:exp-classifcation-attentions}. Results for other setups are additionally provided in Appendix \ref{app:results-cls}}
  \label{tab:classification-config-sweep}
\end{table}

\paragraph{Hyperparameters.}

We use the same architecture as in the induction head experiment. However, we do not use causal masking or positional encodings as the input is effectively a set. We experiment with different numbers of layers and heads, as well as active and inactive MLPs and LayerNorm layers. Full details are provided in Table \ref{tab:classification-config-sweep}.

\paragraph{Metrics.}

We calculate the out-of-distribution flip rate, presented in Figure \ref{fig:main-exp-classification}, as follows. We repeatedly pick two random sequences with different labels from the evaluation set and run a forward pass on them, extracting the attention weights with the \texttt{CLS} token as a query. Then we swap the max-logit tokens for the head $(L, H)$ and check if the model changed the prediction. In pseudocode,
{\small \begin{verbatim}
    X1, X2, Y1, Y2 = pick_samples_with_different_labels()
    
    A1 = attention_matrix(X1)  # [layer, head, query_seq_len, key_seq_len]
    A2 = attention_matrix(X2)

    # obtaining the max-logit tokens excluding CLS itself
    ranking_1[L, H] = A1[L, H,-1,:-1].argsort(reverse=True)
    ranking_2[L, H] = A2[L, H,-1,:-1].argsort(reverse=True)

    # max-logit token from label 2 + others from label 1
    adversarial_1[L, H] = [ranking_2[L, H][0]] + ranking_1[L, H][1:]

    # max-logit token from label 1 + others from label 2
    adversarial_2[L, H] = [ranking_1[L, H][0]] + ranking_2[L, H][1:]

    prediction_adv_1[L, H] = model(adversarial_1[L, H])
    prediction_adv_2[L, H] = model(adversarial_2[L, H])

    # flipped if swapping one token changes the prediction
    flipped_1[L, H] = (prediction_adv_1[L, H] == Y2)
    flipped_2[L, H] = (prediction_adv_2[L, H] == Y1)
    flip_rate[L, H] = (flip_1 + flip_2) / 2
\end{verbatim}}

Figure \ref{fig:cls_task_format} illustrates this procedure. We report \emph{the maximum flip rate} over all heads and layers.

\section{Additional Experimental Results}

\subsection{Value-Softmax Model}

\begin{figure}[h]
  \centering
  \includegraphics[width=1.0\textwidth]{figs/kl.pdf}
  \caption{Training dynamics of Value-Softmax model with KL-divergence as loss function.}
  \label{fig:kl}
\end{figure}

\begin{figure}[h]
  \centering
  \includegraphics[width=1.0\textwidth]{figs/massive_activations_2.pdf}
  \caption{Training dynamics of tied Value-Softmax model ($U \sigma(Ua)$) with logistic loss for different dimensionality $D$. We denote the column of $U$ corresponding to the largest element of $a$ as $U^{(m)}$. In all cases, a sparse solution is found, and the norm of $U^{(m)}$ diverges.} 
  \label{fig:massive_activations}
\end{figure}

\begin{figure}[h]
  \centering
  \includegraphics[width=1.0\textwidth]{figs/condition_numbers.pdf}
  \caption{Training dynamics in regression problems with condition number $\kappa$ from 1 to 5. Higher condition number leads to slower optimization and stronger sparsity.}
  \label{fig:condition_numbers}
\end{figure}

\begin{figure}[t]
  \centering
  \begin{subfigure}[t]{0.46\textwidth}
    \centering
    \includegraphics[width=\textwidth]{figs/square_loss_comparison.pdf}
    \caption{Square loss.}
    %\label{fig:condition_numbers_left}
  \end{subfigure}
  \hfill
  \begin{subfigure}[t]{0.46\textwidth}
    \centering
    \includegraphics[width=\textwidth]{figs/log_loss_comparison.pdf}
    \caption{Logistic loss.}
    %\label{fig:condition_numbers_right}
  \end{subfigure}

  \caption{Training dynamics of value-softmax model for square and logistic losses and various activation functions. Note that square loss does not give rise to sparse solutions. Additionally, sum and sigmoid do not converge to sparse solutions for logistic loss as well.}
  \label{fig:toy-many-plots}
\end{figure}


\FloatBarrier
\subsection{Induction heads acting as attention sinks}
\label{app:results-induction}

% preamble:
% \usepackage{graphicx}
% \usepackage{subcaption}

% ---------- aux = 0.0 ----------
\begin{figure*}[h]
  \centering
  % columns: noshift | shift
  % rows: add, concat, no

  \begin{subfigure}[t]{0.46\textwidth}
    \centering
    \includegraphics[width=\textwidth]{figs/sink_proportions_aux0.0_noshift_add.pdf}
    \caption{additive, $t_i = t_i'$}
  \end{subfigure}\hfill
  \begin{subfigure}[t]{0.46\textwidth}
    \centering
    \includegraphics[width=\textwidth]{figs/sink_proportions_aux0.0_shift_add.pdf}
    \caption{additive, $t_i \ne t_i'$}
  \end{subfigure}

  \vspace{2mm}

  \begin{subfigure}[t]{0.46\textwidth}
    \centering
    \includegraphics[width=\textwidth]{figs/sink_proportions_aux0.0_noshift_concat.pdf}
    \caption{concatenated, $t_i = t_i'$}
  \end{subfigure}\hfill
  \begin{subfigure}[t]{0.46\textwidth}
    \centering
    \includegraphics[width=\textwidth]{figs/sink_proportions_aux0.0_shift_concat.pdf}
    \caption{concatenated, $t_i \ne t_i'$}
  \end{subfigure}

  \vspace{2mm}

  \begin{subfigure}[t]{0.46\textwidth}
    \centering
    \includegraphics[width=\textwidth]{figs/sink_proportions_aux0.0_noshift_no.pdf}
    \caption{no, $t_i = t_i'$}
  \end{subfigure}\hfill
  \begin{subfigure}[t]{0.46\textwidth}
    \centering
    \includegraphics[width=\textwidth]{figs/sink_proportions_aux0.0_shift_no.pdf}
    \caption{no, $t_i \ne t_i'$}
  \end{subfigure}

  \caption{Sink proportions in induction-trained Transformer with $\lambda = 0$. Each row represents a type of positional encodings. Setups that result in training failures are marked by red crosses. We observe that under this setup, most heads duplicate each other and do not form sinks.}
  \label{fig:sink-proportions-aux0}
\end{figure*}

% ---------- aux = 0.2 ----------
\begin{figure*}[h]
  \centering
  % columns: noshift | shift
  % rows: add, concat, no

  \begin{subfigure}[t]{0.46\textwidth}
    \centering
    \includegraphics[width=\textwidth]{figs/sink_proportions_aux0.2_noshift_add.pdf}
    \caption{additive, $t_i = t_i'$}
  \end{subfigure}\hfill
  \begin{subfigure}[t]{0.46\textwidth}
    \centering
    \includegraphics[width=\textwidth]{figs/sink_proportions_aux0.2_shift_add.pdf}
    \caption{additive, $t_i \ne t_i'$}
  \end{subfigure}

  \vspace{2mm}

  \begin{subfigure}[t]{0.46\textwidth}
    \centering
    \includegraphics[width=\textwidth]{figs/sink_proportions_aux0.2_noshift_concat.pdf}
    \caption{concatenated, $t_i = t_i'$}
  \end{subfigure}\hfill
  \begin{subfigure}[t]{0.46\textwidth}
    \centering
    \includegraphics[width=\textwidth]{figs/sink_proportions_aux0.2_shift_concat.pdf}
    \caption{concatenated, $t_i \ne t_i'$}
  \end{subfigure}

  \vspace{2mm}

  \begin{subfigure}[t]{0.46\textwidth}
    \centering
    \includegraphics[width=\textwidth]{figs/sink_proportions_aux0.2_noshift_no.pdf}
    \caption{no, $t_i = t_i'$}
  \end{subfigure}\hfill
  \begin{subfigure}[t]{0.46\textwidth}
    \centering
    \includegraphics[width=\textwidth]{figs/sink_proportions_aux0.2_shift_no.pdf}
    \caption{no, $t_i \ne t_i'$}
  \end{subfigure}

  \caption{Sink proportions in induction-trained Transformer with $\lambda = 0.2$. Each row represents a type of positional encodings. Setups that result in training failures are marked by red crosses. Interestingly, we observe that setups with no explicit positional encodings lead to an increased formation of attention sinks. Note that in this case the model has access to positional information through causal masking, so the task is still solvable.}
  \label{fig:sink-proportions-aux02}
\end{figure*}

% ---------- accuracy, aux = 0.0 ----------
\begin{figure*}[h]
  \centering
  % columns: noshift | shift
  % rows: add, concat, no

  \begin{subfigure}[t]{0.46\textwidth}
    \centering
    \includegraphics[width=\textwidth]{figs/accuracy_aux0.0_noshift_add.pdf}
    \caption{additive, $t_i = t_i'$}
  \end{subfigure}\hfill
  \begin{subfigure}[t]{0.46\textwidth}
    \centering
    \includegraphics[width=\textwidth]{figs/accuracy_aux0.0_shift_add.pdf}
    \caption{additive, $t_i \ne t_i'$}
  \end{subfigure}

  \vspace{2mm}

  \begin{subfigure}[t]{0.46\textwidth}
    \centering
    \includegraphics[width=\textwidth]{figs/accuracy_aux0.0_noshift_concat.pdf}
    \caption{concatenated, $t_i = t_i'$}
  \end{subfigure}\hfill
  \begin{subfigure}[t]{0.46\textwidth}
    \centering
    \includegraphics[width=\textwidth]{figs/accuracy_aux0.0_shift_concat.pdf}
    \caption{concatenated, $t_i \ne t_i'$}
  \end{subfigure}

  \vspace{2mm}

  \begin{subfigure}[t]{0.46\textwidth}
    \centering
    \includegraphics[width=\textwidth]{figs/accuracy_aux0.0_noshift_no.pdf}
    \caption{no, $t_i = t_i'$}
  \end{subfigure}\hfill
  \begin{subfigure}[t]{0.46\textwidth}
    \centering
    \includegraphics[width=\textwidth]{figs/accuracy_aux0.0_shift_no.pdf}
    \caption{no, $t_i \ne t_i'$}
  \end{subfigure}

  \caption{Average induction accuracy for models trained with $\lambda = 0$. Each row represents a type of positional encodings.}
  \label{fig:accuracy-aux0}
\end{figure*}

% ---------- accuracy, aux = 0.2 ----------
\begin{figure*}[h]
  \centering
  % columns: noshift | shift
  % rows: add, concat, no

  \begin{subfigure}[t]{0.46\textwidth}
    \centering
    \includegraphics[width=\textwidth]{figs/accuracy_aux0.2_noshift_add.pdf}
    \caption{additive, $t_i = t_i'$}
  \end{subfigure}\hfill
  \begin{subfigure}[t]{0.46\textwidth}
    \centering
    \includegraphics[width=\textwidth]{figs/accuracy_aux0.2_shift_add.pdf}
    \caption{additive, $t_i \ne t_i'$}
  \end{subfigure}

  \vspace{2mm}

  \begin{subfigure}[t]{0.46\textwidth}
    \centering
    \includegraphics[width=\textwidth]{figs/accuracy_aux0.2_noshift_concat.pdf}
    \caption{concatenated, $t_i = t_i'$}
  \end{subfigure}\hfill
  \begin{subfigure}[t]{0.46\textwidth}
    \centering
    \includegraphics[width=\textwidth]{figs/accuracy_aux0.2_shift_concat.pdf}
    \caption{concatenated, $t_i \ne t_i'$}
  \end{subfigure}

  \vspace{2mm}

  \begin{subfigure}[t]{0.46\textwidth}
    \centering
    \includegraphics[width=\textwidth]{figs/accuracy_aux0.2_noshift_no.pdf}
    \caption{no, $t_i = t_i'$}
  \end{subfigure}\hfill
  \begin{subfigure}[t]{0.46\textwidth}
    \centering
    \includegraphics[width=\textwidth]{figs/accuracy_aux0.2_shift_no.pdf}
    \caption{no, $t_i \ne t_i'$}
  \end{subfigure}

  \caption{Average induction accuracy for models trained with $\lambda = 0.2$. Each row represents a type of positional encodings.}
  \label{fig:accuracy-aux02}
\end{figure*}

\FloatBarrier
\subsection{Emerging imbalance in token influence on classifier's predictions}
\label{app:results-cls}

% preamble:
% \usepackage{graphicx}
% \usepackage{subcaption}

% ---------------- vocab = 41 ----------------
\begin{figure*}[h]
  \centering

  \begin{subfigure}[t]{0.46\textwidth}
    \centering
    \includegraphics[width=\textwidth]{figs/fliprate_vocab41_mlp_ln.pdf}
    \caption{MLP, LN}
  \end{subfigure}\hfill
  \begin{subfigure}[t]{0.46\textwidth}
    \centering
    \includegraphics[width=\textwidth]{figs/fliprate_vocab41_mlp_noln.pdf}
    \caption{MLP, no LN}
  \end{subfigure}

  \vspace{2mm}

  \begin{subfigure}[t]{0.46\textwidth}
    \centering
    \includegraphics[width=\textwidth]{figs/fliprate_vocab41_nomlp_ln.pdf}
    \caption{no MLP, LN}
  \end{subfigure}\hfill
  \begin{subfigure}[t]{0.46\textwidth}
    \centering
    \includegraphics[width=\textwidth]{figs/fliprate_vocab41_nomlp_noln.pdf}
    \caption{no MLP, no LN}
  \end{subfigure}

  \caption{Flip-rate plots for vocabulary size $|\mathcal{V}|=41$. Rows: MLP vs. no-MLP. Columns: LayerNorm vs. no-LayerNorm.}
  \label{fig:fliprate-vocab41}
\end{figure*}

% ---------------- vocab = 81 ----------------
\begin{figure*}[h]
  \centering

  \begin{subfigure}[t]{0.46\textwidth}
    \centering
    \includegraphics[width=\textwidth]{figs/fliprate_vocab81_mlp_ln.pdf}
    \caption{MLP, LN}
  \end{subfigure}\hfill
  \begin{subfigure}[t]{0.46\textwidth}
    \centering
    \includegraphics[width=\textwidth]{figs/fliprate_vocab81_mlp_noln.pdf}
    \caption{MLP, no LN}
  \end{subfigure}

  \vspace{2mm}

  \begin{subfigure}[t]{0.46\textwidth}
    \centering
    \includegraphics[width=\textwidth]{figs/fliprate_vocab81_nomlp_ln.pdf}
    \caption{no MLP, LN}
  \end{subfigure}\hfill
  \begin{subfigure}[t]{0.46\textwidth}
    \centering
    \includegraphics[width=\textwidth]{figs/fliprate_vocab81_nomlp_noln.pdf}
    \caption{no MLP, no LN}
  \end{subfigure}

  \caption{Flip-rate plots for vocabulary size $|\mathcal{V}|=81$. Rows: MLP vs. no-MLP. Columns: LayerNorm vs. no-LayerNorm.}
  \label{fig:fliprate-vocab81}
\end{figure*}

% ---------------- vocab = 161 ----------------
\begin{figure*}[h]
  \centering

  \begin{subfigure}[t]{0.46\textwidth}
    \centering
    \includegraphics[width=\textwidth]{figs/fliprate_vocab161_mlp_ln.pdf}
    \caption{MLP, LN}
  \end{subfigure}\hfill
  \begin{subfigure}[t]{0.46\textwidth}
    \centering
    \includegraphics[width=\textwidth]{figs/fliprate_vocab161_mlp_noln.pdf}
    \caption{MLP, no LN}
  \end{subfigure}

  \vspace{2mm}

  \begin{subfigure}[t]{0.46\textwidth}
    \centering
    \includegraphics[width=\textwidth]{figs/fliprate_vocab161_nomlp_ln.pdf}
    \caption{no MLP, LN}
  \end{subfigure}\hfill
  \begin{subfigure}[t]{0.46\textwidth}
    \centering
    \includegraphics[width=\textwidth]{figs/fliprate_vocab161_nomlp_noln.pdf}
    \caption{no MLP, no LN}
  \end{subfigure}

  \caption{Flip-rate plots for vocabulary size $|\mathcal{V}|=161$. Rows: MLP vs. no-MLP. Columns: LayerNorm vs. no-LayerNorm.}
  \label{fig:fliprate-vocab161}
\end{figure*}

% preamble:
% \usepackage{graphicx}
% \usepackage{subcaption}

\begin{figure*}[t]
  \centering
  % 3 rows: heads = 1,2,4
  % 2 cols: layers = 1,2
  % fixed: vocab41, nomlp, noln

  % -------- row: heads = 1 --------
  \begin{subfigure}[t]{0.46\textwidth}
    \centering
    \includegraphics[width=\textwidth]{figs/fliprate_vocab41_layers1_heads1_nomlp_noln.pdf}
    \caption{$H{=}1$, $L{=}1$}
  \end{subfigure}\hfill
  \begin{subfigure}[t]{0.46\textwidth}
    \centering
    \includegraphics[width=\textwidth]{figs/fliprate_vocab41_layers2_heads1_nomlp_noln.pdf}
    \caption{$H{=}1$, $L{=}2$}
  \end{subfigure}

  \vspace{2mm}

  % -------- row: heads = 2 --------
  \begin{subfigure}[t]{0.46\textwidth}
    \centering
    \includegraphics[width=\textwidth]{figs/fliprate_vocab41_layers1_heads2_nomlp_noln.pdf}
    \caption{$H{=}2$, $L{=}1$}
  \end{subfigure}\hfill
  \begin{subfigure}[t]{0.46\textwidth}
    \centering
    \includegraphics[width=\textwidth]{figs/fliprate_vocab41_layers2_heads2_nomlp_noln.pdf}
    \caption{$H{=}2$, $L{=}2$}
  \end{subfigure}

  \vspace{2mm}

  % -------- row: heads = 4 --------
  \begin{subfigure}[t]{0.46\textwidth}
    \centering
    \includegraphics[width=\textwidth]{figs/fliprate_vocab41_layers1_heads4_nomlp_noln.pdf}
    \caption{$H{=}4$, $L{=}1$}
  \end{subfigure}\hfill
  \begin{subfigure}[t]{0.46\textwidth}
    \centering
    \includegraphics[width=\textwidth]{figs/fliprate_vocab41_layers2_heads4_nomlp_noln.pdf}
    \caption{$H{=}4$, $L{=}2$}
  \end{subfigure}

  \caption{Flip-rate plots for $|\mathcal{V}|=41$ with no MLP and no LayerNorm. Rows: number of heads $H\in\{1,2,4\}$. Columns: number of layers $L\in\{1,2\}$. Increasing the number of heads or layers leads to the diminishing flip rate.}
  \label{fig:fliprate-vocab41-3x2-heads-layers-nomlp-noln}
\end{figure*}


\end{document}

% This document was modified from the file originally made available by
% Pat Langley and Andrea Danyluk for ICML-2K. This version was created
% by Iain Murray in 2018, and modified by Alexandre Bouchard in
% 2019 and 2021 and by Csaba Szepesvari, Gang Niu and Sivan Sabato in 2022.
% Modified again in 2023 and 2024 by Sivan Sabato and Jonathan Scarlett.
% Previous contributors include Dan Roy, Lise Getoor and Tobias
% Scheffer, which was slightly modified from the 2010 version by
% Thorsten Joachims & Johannes Fuernkranz, slightly modified from the
% 2009 version by Kiri Wagstaff and Sam Roweis's 2008 version, which is
% slightly modified from Prasad Tadepalli's 2007 version which is a
% lightly changed version of the previous year's version by Andrew
% Moore, which was in turn edited from those of Kristian Kersting and
% Codrina Lauth. Alex Smola contributed to the algorithmic style files.
